\makeatletter
\def\input@path{{../}{./}}
\makeatother
\documentclass[10pt,a4paper,oneside]{ctexbook}
\RequirePackage{xcolor}
\definecolor{main}{RGB}{0,82,136}
\definecolor{light}{RGB}{230,240,250}

\usepackage{hxnotebook}

\begin{document}

% cover
\makecover
  {环境卫生学}       
  {华西大纲·PPT·历年真题整理}    
  {复习笔记}                  
  {\today}                 
  {greedyCat}                  

% ===== Table of Contents =====
\tableofcontents

% ===== 使用说明 =====
\chapter*{使用说明}
\addcontentsline{toc}{chapter}{使用说明}

\begin{itemize}[leftmargin=*]
    \item 本复习笔记严格依照《环境卫生学》教学大纲（2026版）章节编排，重点内容与大纲标注的掌握内容一致。
    \item 各章末附有复习题（名词解释、简答题、非标准答案题），题目主要来源于教学平台题库及大纲重点内容。
    \item \textbf{\color{keyred}红色框}标识的内容为必须重点掌握的核心知识。
    \item \textbf{\color{main}主题框}为概念释义，帮助理解专业术语。
    \item 建议结合教师授课PPT中的标红内容重点复习。
\end{itemize}

\vspace{1cm}
\begin{center}
\begin{tcolorbox}[colback=yellow!10, colframe=yellow!80!black, boxrule=0.8pt, arc=6pt, width=0.85\textwidth, breakable]
\textbf{考核形式提示：}
\begin{itemize}[leftmargin=*]
    \item 平时考核成绩 50\% + 期末理论考核 50\%
    \item 考试范围：以教学大纲和教师课堂讲授内容为主
    \item 大纲下划线内容 = 重点掌握内容（本笔记已用红色框标注）
\end{itemize}
\end{tcolorbox}
\end{center}

\newpage

% =====================================================================
%  CHAPTER 1: 绪论
% =====================================================================
\chapter{绪论}
\setcounter{chapter}{1}
\chapterintro{
掌握环境卫生学的定义、研究对象和内容；了解我国环境卫生工作主要成就及全球环境问题；理解"绿水青山就是金山银山"的发展理念。
}

% ----- 1.1 概述 -----
\section{环境卫生学概述}

\begin{defbox}
\textbf{环境卫生学（Environmental Hygiene / Environmental Health）}是研究自然环境和生活环境与人群健康的关系，揭示环境因素对人群健康影响的发生、发展规律，为充分利用环境有益因素和控制环境有害因素提出卫生要求和预防对策，增进人体健康，维护和提高人体健康水平的学科。
\end{defbox}

\subsection{研究对象与内容}

环境卫生学以\keyword{人类及其周围的环境}为研究对象，包括自然环境与生活环境。

\begin{tcolorbox}[colback=light, colframe=main, boxrule=0.5pt, arc=4pt, title=\textbf{环境卫生学的研究内容}, breakable]
\begin{enumerate}[leftmargin=*, itemsep=3pt]
    \item \imp{环境与健康关系的基础理论研究}——揭示环境因素对健康的影响机制。包括：
    \begin{itemize}
        \item \textbf{人类基因组计划（HGP）}：完成了人类23对染色体约60亿个核苷酸排列顺序的测定
        \item \textbf{环境基因组计划（EGP）}：研究有重要功能的环境应答基因多态性，确定引起环境暴露致病危险性差异的遗传因素
        \item \textbf{国际人类基因组单体型图谱（HapMap）}：构建不同人群的高密度单核苷酸多态（SNP）图谱，揭示群体分子遗传机制
    \end{itemize}
    \item \imp{环境因素与健康关系的确认性研究}——运用流行病学和毒理学方法。环境因素对健康影响的特点：\textbf{长期、低剂量、反复作用}
    \item \imp{创建和引进环境卫生学研究的新技术和新方法}——提升研究精度与深度
    \item \imp{研究环境卫生监督体系的理论依据}——为卫生标准与法规提供科学支撑
\end{enumerate}
\end{tcolorbox}

\subsection{环境的基本概念}

\begin{defbox}
\textbf{环境（Environment）}：指作用于人类所有外界力量（因素）的总和。人与环境既相互对立，又相互制约，既相互依存又相互转化，两者存在对立统一关系。
\end{defbox}

\begin{defbox}
\textbf{环境介质（Environmental Media）}：人类赖以生存的物质环境条件，通常以气态、液态和固态三种物质形态存在，能够容纳和运载各种环境因素。主要指\keyword{大气、水、土壤（岩石）以及包括人体在内的所有生物体}。
\end{defbox}

\textbf{环境介质的特点：}
\begin{enumerate}[leftmargin=*]
    \item 三种物质形态往往不完全以单一介质形式存在
    \item 三种物质形态在一定条件下可以相互转化
    \item 环境介质的运动可携带环境污染物向远方扩散
    \item 环境介质还能维持自身的稳定状态
\end{enumerate}

\begin{defbox}
\textbf{环境因素（Environmental Factors）}：被介质容纳和转运的成分或介质中各种无机和有机的组成成分。分为\keyword{物理性、化学性、生物性}三类。

\begin{itemize}[leftmargin=*]
    \item \textbf{物理因素：}小气候（气温、气湿、气流和热辐射，对机体热平衡产生明显影响）、噪声、振动、非电离辐射、电离辐射
    \item \textbf{化学因素：}无机和有机化学物
    \item \textbf{生物因素：}细菌、真菌、病毒、寄生虫和生物性变应原（如植物花粉、真菌孢子等）
\end{itemize}
\end{defbox}

\subsection{原生环境与次生环境}

\begin{defbox}
\textbf{原生环境（Primary Environment）}：天然形成的未受或少受人为因素影响的环境。有利方面：清洁的空气和水、充足的日光、适宜的小气候、秀丽的自然风光。不利方面：自然灾害、天然植物毒素、天然动物毒素、\textbf{生物地球化学性疾病}。

\textbf{次生环境（Secondary Environment）}：受人为活动影响形成的环境。有利方面：水利工程、资源开采、化工行业发展等。不利方面：污染环境，影响健康——全球气候变暖、臭氧层破坏、酸雨、生物多样性锐减；\textbf{环境污染性疾病}。
\end{defbox}

\begin{defbox}
\textbf{环境内分泌干扰物（EDCs）}：许多环境化学污染物（如有机氯化合物、二噁英、烷基酚、邻苯二甲酸酯等）对维持机体内环境稳态和调节发育过程的体内天然激素的生成、释放、转运、代谢、结合、效应造成严重影响的外源性物质，被称为内分泌干扰化学物。
\end{defbox}

% ----- 1.2 污染物 -----
\section{污染物分类}

\begin{defbox}
\textbf{一次污染物（Primary Pollutant）}：由污染源直接排入环境未发生变化的污染物。
\end{defbox}

\begin{defbox}
\textbf{二次污染物（Secondary Pollutant）}：一次污染物进入环境后在物理、化学或生物学作用下，或与其他物质发生反应而形成与原来污染物的理化性质和毒性完全不同的新污染物。
\end{defbox}

\begin{keybox}
常见的二次污染物：光化学烟雾中的臭氧、过氧乙酰硝酸酯（PAN）、酸雨中的硫酸和硝酸等。二次污染物的毒性往往比一次污染物更大。
\end{keybox}

\begin{defbox}
\textbf{新污染物（Emerging Pollutants）}：由人类活动造成的、目前已在环境中明确存在、危害人体健康和生态环境的，但因其生产使用历史相对较短或危害发现较晚，尚无法律法规和标准予以明确规定或规定不完善的所有在生产建设或其他活动中产生的污染物。一般包括环境激素、抗生素等新化学物质和微/纳塑料等细颗粒物。

\textbf{环境自净（Environmental Self-Purification）}：环境受到污染后，在物理、化学和生物的作用下，对有害物质进行扩散、稀释、氧化还原、生物降解等作用，逐步降低污染物的浓度，减小甚至消除污染物的毒性。

\textbf{公害病（Public Nuisance Disease）}：严重的地区性环境污染性疾病，由政府认定，须经严格鉴定和国家法律的正式认可，有\imp{医学和法律的双重含义}，病人享受国家补贴。

\textbf{毒物兴奋效应（Hormesis Effect）}：某些物质在低剂量时对生物系统具有刺激（有益）作用，而在高剂量时具有抑制（有害）作用，其剂量-效应关系以\imp{双向曲线（倒U或J形）}为特征。典型环境污染物如镉、铅、汞、二噁英等均表现出Hormesis现象。环境因素对人体健康的这种\imp{双重性}（Dual Role）是环境卫生学的重要基础概念。
\end{defbox}

% ----- 1.3 环境问题 -----
\section{全球环境问题与挑战}

\subsection{全球性环境问题}
\begin{enumerate}[leftmargin=*]
    \item \keyword{全球气候变暖}
    \item \keyword{臭氧层破坏}
    \item \keyword{酸雨}
    \item \keyword{生物多样性锐减}
\end{enumerate}

世界第一部关于环境的著作：\imp{《寂静的春天》}（Rachel Carson，1962），标志着现代环保运动的开端。

\textbf{我国环境保护工作的发展历程：}充分理解——没有安全的环境，就不可能有安全的饮水、食品、空气，更不可能有人类健康。

\subsection{世界八大公害事件}

\begin{table}[htbp]
\centering
\caption{世界八大公害事件简表}
\begin{tabularx}{\textwidth}{lXXX}
\hline
\textbf{事件名称} & \textbf{时间} & \textbf{污染物} & \textbf{健康危害} \\
\hline
马斯河谷烟雾事件 & 1930（比利时） & SO$_2$、粉尘 & 呼吸系统疾病 \\
多诺拉烟雾事件 & 1948（美国） & SO$_2$、金属粉尘 & 呼吸系统疾病 \\
伦敦烟雾事件 & 1952（英国） & SO$_2$、烟尘 & 急性中毒、死亡 \\
洛杉矶光化学烟雾 & 1943起（美国） & O$_3$、PAN & 眼睛刺激、呼吸疾病 \\
水俣病 & 1953起（日本） & 甲基汞 & 神经系统损害 \\
痛痛病 & 1955起（日本） & 镉 & 骨骼病变 \\
四日市哮喘 & 1961起（日本） & SO$_2$、粉尘 & 哮喘 \\
米糠油事件 & 1968（日本） & 多氯联苯 & 皮肤病变 \\
\hline
\end{tabularx}
\end{table}

\begin{keybox}
\textbf{水俣病}（Minamata Disease）：由甲基汞污染水体，通过食物链生物放大作用引起慢性甲基汞中毒，主要表现为神经系统损害。

\textbf{痛痛病}（Itai-itai Disease）：由镉污染土壤和水体，通过食物链蓄积引起慢性镉中毒，以骨骼病变和肾损害为特征。
\end{keybox}

% ----- 1.4 复习题 -----
\section{复习题}

\begin{enumerate}[leftmargin=*, itemsep=8pt]
    \item \textbf{【名词解释】}环境卫生学；新污染物（Emerging Pollutants）
    \item \textbf{【名词解释】}一次污染物与二次污染物；环境自净
    \item \textbf{【名词解释】}原生环境与次生环境；公害病
    \item \textbf{【名词解释】}Hormesis效应（毒物兴奋效应）；环境内分泌干扰物（EDCs）
    \item \textbf{【简答题】}简述环境卫生学的研究内容。
    \item \textbf{【简答题】}简述我国环境卫生工作的主要成就及今后任务。
    \item \textbf{【非标准答案题】}试述你对"环境就是民生，青山就是美丽，蓝天也是幸福"的理解，以及环境卫生学在"美丽中国"建设中应发挥的作用。
\end{enumerate}

\subsection*{参考答案要点}

\begin{enumerate}[leftmargin=*, itemsep=5pt]
    \item \textbf{环境卫生学}：研究自然环境和生活环境与人群健康的关系，揭示环境因素对人群健康影响的发生、发展规律，为充分利用环境有益因素和控制环境有害因素提出卫生要求和预防对策，增进人体健康，维护和提高人体健康水平的学科。\textbf{新污染物}：由人类活动造成的、在环境中明确存在但尚无法律法规和标准予以明确规定的污染物，如环境激素、抗生素、微/纳塑料等。

    \item \textbf{一次污染物}：由污染源直接排入环境未发生变化的污染物。\textbf{二次污染物}：一次污染物进入环境后在物理、化学或生物学作用下转化形成的新污染物，其理化性质和毒性通常与原污染物不同。\textbf{环境自净}：环境受污染后，在物理、化学和生物作用下逐步降低污染物浓度、减小或消除污染物毒性的过程。

    \item \textbf{原生环境}：天然形成的未受或少受人为因素影响的环境。\textbf{次生环境}：受人为活动影响形成的环境。\textbf{公害病}：严重的地区性环境污染性疾病，有医学和法律的双重含义。

    \item 环境卫生学的四大研究内容：①环境与健康关系的基础理论研究；②环境因素与健康关系的确认性研究；③创建和引进环境卫生学研究的新技术和新方法；④研究环境卫生监督体系的理论依据。

    \item 我国环境卫生工作的主要成就包括：建立了较为完善的环境卫生标准体系；开展了广泛的环境污染与健康关系的调查研究；积极参与全球环境治理。今后任务：加强基础理论和确认性研究；完善环境卫生法律法规标准体系；加强农村环境卫生工作；开拓新污染物、气候变化与健康等新领域。

    \item （开放性问题，围绕环境保护与人民健康的关系、可持续发展理念、公共卫生工作者的责任等方面展开论述。）
\end{enumerate}

\newpage

% =====================================================================
%  CHAPTER 2: 环境与健康的关系
% =====================================================================
\chapter{环境与健康的关系}
\setcounter{chapter}{2}
\chapterintro{
本章是环境卫生学的理论基础，共12学时（占比最大）。掌握生态系统、食物链、生物放大等基本概念；掌握人群健康效应谱、剂量-反应关系；掌握环境流行病学和毒理学研究方法的特点；掌握健康危险度评价的四个阶段。
}

% ----- 2.1 人类的环境 -----
\section{人类的环境}

\subsection{生态系统}
\begin{defbox}
\textbf{生物圈（Biosphere）}：从海平面以下深约12 km至海平面以上高约10 km的范围，包括一部分大气圈、水圈和土壤岩石圈，是地球上所有生命物质及其生存环境的整体。
\end{defbox}

\begin{defbox}
\textbf{生态系统（Ecosystem）}：在一定空间范围内，由\keyword{生物群落及其环境}组成，借助各种功能流（\imp{物质流、能量流、物种流和信息流}）所联结的稳态系统。
\end{defbox}

\textbf{生态系统组成：}
\begin{itemize}[leftmargin=*]
    \item \imp{生产者}（绿色植物、光合细菌）——将无机物转化为有机物
    \item \imp{消费者}（动物）——直接或间接以生产者为食
    \item \imp{分解者}（微生物）——将有机物分解为无机物
    \item \imp{无机界}（阳光、水、空气、土壤等）
\end{itemize}

\textbf{生态系统的主要特征：}
\begin{enumerate}[leftmargin=*]
    \item 整体性——各要素对系统整体性起作用
    \item 开放性——与外界进行物质交换
    \item 自调控——种群密度调控、种间数量调控、生物与环境适应调控
    \item 可持续性——不断进行物质循环和能量流动
\end{enumerate}

\begin{defbox}
\textbf{生态系统健康（Ecosystem Health）}：具有活力、结构稳定和自调节能力的生态系统，是生态系统的综合特性。
\end{defbox}

\subsection{食物链与生物放大}

\begin{defbox}
\textbf{食物链（Food Chain）}：生态系统中不同营养级的生物逐级被吞食以满足生存需要而建立起来的链锁关系。
\end{defbox}

\begin{defbox}
\textbf{生物富集作用（Bioenrichment）}：生物从环境中摄入浓度极低的重金属元素或难降解化合物，在体内逐渐累积，使生物体内该物质的浓度大大超过环境中的浓度。

\textbf{浓集系数} = 污染物在生物体内浓度 / 污染物在环境中浓度。
\end{defbox}

\begin{defbox}
\textbf{迁移（Transportation）}：环境物质在环境中发生的空间位移过程，伴随污染物在环境介质中浓度改变。

\textbf{转化（Transform）}：化学物（或污染物）在环境中通过物理、化学或生物作用改变其存在形态或转变成另一物质的过程。
\end{defbox}

\begin{keybox}
\textbf{环境污染物的迁移和转化对环境因素暴露的影响：}
\begin{enumerate}[leftmargin=*]
    \item 扩大暴露范围
    \item 增加暴露途径
    \item 影响暴露剂量
    \item 改变污染物性质和毒性
\end{enumerate}

\textbf{生物放大作用（Biomagnification）}：环境中的重金属元素和难降解有机物，通过食物链转移到高位营养级生物体内，使其浓度逐级在生物体内放大。

\textbf{典型实例：}DDT在水生食物链中的迁移——浮游生物（0.04 ppm）→ 小鱼（0.5 ppm）→ 大鱼（2.0 ppm）→ 水鸟（25 ppm），浓缩系数达600倍以上。

\textbf{水俣病}和\textbf{痛痛病}就是生物放大作用的经典案例。
\end{keybox}

% ----- 2.2 人与环境的辩证关系 -----
\section{人与环境的辩证统一关系}

\subsection{人与环境在物质上的统一性}
人体血液中的60多种元素含量与地壳中这些元素的含量呈明显正相关，说明人体组成与地球化学环境具有\keyword{物质上的统一性}。

\subsection{人类对环境的适应性}
\begin{itemize}[leftmargin=*]
    \item 光适应
    \item 高原低氧适应
    \item 热适应
    \item 人类进化适应
\end{itemize}

\textbf{生物学基础：}环境应答基因（Environmental Response Gene）的多态性是造成人群易感性差异的重要原因。涉及\imp{基因多态性}和\imp{单核苷酸多态性（SNP）}。

\begin{defbox}
\textbf{环境基因交互作用}：指个体基因与外部环境相互作用的过程。①环境影响基因的表达——如某些环境化学物既是CYP酶的底物又是酶的诱导剂，可增强CYP酶的代谢能力，从而增毒或解毒；②基因多态性影响有害因素的作用方式——如影响化学物在体内的吸收、转运、分布、代谢等。

典型实例：\textbf{苯丙酮尿症}——苯丙氨酸羟化酶缺乏且摄入苯丙氨酸；\textbf{肿瘤发生}——致癌因素暴露 + DNA修复基因SNP + 代谢酶遗传多态性。

机体适应分为\imp{群体适应}（通过基因或表观遗传改变实现，可代际遗传）和\imp{个体适应}（不涉及遗传物质改变，如个体对高温、高原的适应）。
\end{defbox}

\subsection{环境因素对健康影响的双重性}
环境因素既有\imp{有益作用}（如紫外线促进维生素D合成），也有\imp{有害作用}（如过量紫外线导致皮肤癌）。

% ----- 2.3 暴露特征 -----
\section{环境因素的暴露特征}

\subsection{暴露的基本概念}

\begin{defbox}
\textbf{暴露（Exposure）}：人体接触某一有害环境因素的过程。
\end{defbox}

\begin{itemize}[leftmargin=*]
    \item \imp{外剂量（External Dose）}：环境介质中某种环境因素的浓度或含量，根据人体接触特征估计个体的暴露水平。
    \item \imp{内剂量（Internal Dose）}：机体内已吸收的污染物的量。
    \item \imp{生物有效剂量（Biologically Effective Dose）}：经吸收、代谢活化、转运最终到达靶部位或替代性靶部位的污染物量。
\end{itemize}

\begin{defbox}
\textbf{生物半减期（Biological Half-life, t$_{1/2}$）}：化合物在体内含量减少一半所需的时间。化合物进入机体经历六个生物半减期后，体内最大可能蓄积量趋于稳定，此后摄入量与排出量趋于平衡。摄入量愈大，达到平衡后其最大蓄积量也愈大；摄入量若减少到致使体内最大蓄积量低于产生有害效应的水平，长期作用也观察不到对机体的有害效应。
\end{defbox}

\subsection{环境多因素暴露与联合作用}

\begin{table}[htbp]
\centering
\caption{联合作用的类型}
\begin{tabularx}{\textwidth}{lX}
\hline
\textbf{类型} & \textbf{定义及举例} \\
\hline
\imp{相加作用（Additive）} & 多种化学物联合作用强度等于各化学物单独作用强度之和。如同系物或毒作用靶器官相同的化学物。大部分刺激性气体的刺激作用、两种有机磷农药抑制胆碱酯酶的作用 \\
\imp{独立作用（Independent）} & 各化学物作用于不同受体/部位/靶器官，引发的生物效应互不干扰，表现为各自毒性效应。如酒精与氯乙烯的联合作用 \\
\imp{协同作用（Synergistic）} & 联合作用强度大于各化学物单独作用强度之和。如马拉硫磷和苯硫磷的联合作用 \\
\imp{增强作用（Potentiation）} & 一种化学物本身对某器官无毒性，但能增强另一种化学物的毒性。如异丙醇对肝脏无毒性，但与四氯化碳同时进入机体时使四氯化碳毒性远高于单独作用 \\
\imp{拮抗作用（Antagonism）} & 联合作用强度小于各化学物单独作用强度之和。包括：
\begin{itemize}
    \item 竞争作用：肟类化合物与有机磷竞争胆碱酯酶
    \item 功能性拮抗：阿托品对抗有机磷的毒蕈碱症状
    \item 代谢过程变化：卤代苯类诱导有机磷代谢使毒性减弱
\end{itemize} \\
\hline
\end{tabularx}
\end{table}

\subsection{剂量-反应关系}

\begin{defbox}
\textbf{无阈值化合物（No-threshold Compound）}：在大于零的剂量暴露下均可能发生有害效应的化合物。其剂量-反应曲线延长线通过坐标原点，\textbf{无安全剂量}。典型代表：遗传毒性致癌物。

\textbf{阈值化合物（Threshold Compound）}：仅在达到或大于阈剂量时才产生效应。剂量-反应曲线多呈\textbf{S形或抛物线形}。有两个阈值的化合物呈\textbf{U形}曲线。
\end{defbox}

\begin{keybox}
\textbf{阈值理论在环境卫生标准制定中的重要性：}对于有阈化合物，可通过毒理学和流行病学方法确定阈剂量，进而除以安全系数得到卫生标准限值。无阈化合物（如遗传毒性致癌物）则不适用阈值理论，需采用健康风险评价方法。
\end{keybox}

\begin{tipbox}
\textbf{研究环境因素联合作用的意义：}
\begin{enumerate}[leftmargin=*]
    \item 多种环境因素同时存在时对人体的作用与其中一种单独存在时有所不同，往往呈现复杂的交互作用，使机体的毒性效应发生改变
    \item 在阐明环境因素对健康的影响、制订混合污染物的卫生标准、开展多种有害因素共同作用的危险度评价、采取防治对策等方面具有实际意义
\end{enumerate}
\end{tipbox}

% ----- 2.4 健康效应 -----
\section{环境因素的健康效应}

\subsection{人群健康效应谱}

\begin{keybox}
\textbf{人群健康效应谱（Spectrum of Health Effect）}：人群暴露于环境污染物时产生的不良反应在性质、程度和范围上表现的征象，从弱到强分为五级：

\begin{enumerate}[leftmargin=*]
    \item \textbf{污染物负荷增加}——体内负荷增加，不引起生理功能和生化代谢变化
    \item \textbf{生理代偿性变化}——出现某些生理功能和生化代谢变化（生理代偿）
    \item \textbf{亚临床状态（准病态）}——生理代谢异常，机体处于病理性代偿调节状态，无明显临床症状
    \item \textbf{临床性疾病}——机体功能失调，出现临床症状
    \item \textbf{死亡}——严重中毒导致死亡
\end{enumerate}

\textbf{分布规律：}呈金字塔形分布——健康人群最多，亚临床状态居中，临床患者较少，死亡最少。
\end{keybox}

\begin{defbox}
\textbf{高危人群（High Risk Group）}：对环境暴露的健康效应出现较早、效应较显著、易感性较高的人群，包括高暴露人群（职业人群和特殊暴露人群）和高敏感人群（易感人群）。

\textbf{影响人群易感性的因素：}
\begin{itemize}[leftmargin=*]
    \item \imp{非遗传因素}：年龄、健康状况、营养状态、生活习惯、暴露史、心理状态、保护措施
    \item \imp{遗传因素}：性别、种族、遗传缺陷、基因多态性
\end{itemize}

\textbf{脆弱人群（Vulnerable Group）}：身体状况较差或抗病能力较弱的人。

\textbf{敏感人群（Sensitive Group）}：接触有害物质时，由于个体生物学因素使其毒性反应的出现较普通人群更早、反应更强的人群。
\end{defbox}

\begin{tipbox}
\textbf{研究高危人群的意义：}
\begin{enumerate}[leftmargin=*]
    \item 是最早出现健康效应的人群，有助于全人群健康效应的早发现和早预防
    \item 高危人群的健康效应更复杂和严重，对该人群的保护即是对全人群的保护
    \item 对高危人群提早筛查可获得最大流行病学和卫生经济学效益
\end{enumerate}
\end{tipbox}

\subsection{环境污染对人群健康的危害}

\begin{enumerate}[leftmargin=*]
    \item \imp{急性危害（Acute Hazard）}：污染物短时间内大量进入环境，使暴露人群短时间内出现不良反应、急性中毒甚至死亡。包括：大气污染烟雾事件、过量排放和事故性排放、生物性污染引起的急性传染病。
    \item \imp{慢性危害（Chronic Hazard）}：有害污染物以低浓度、长时间反复作用于机体产生的危害。包括非特异性影响、慢性疾患（COPD等）、持续性蓄积危害。
    \item \imp{致癌致畸危害}：如反应停（Thalidomide）事件——孕妇服用导致胎儿畸形。
    \item \imp{环境致畸物和致畸因素}：孕期摄入环境污染物可引起胎儿畸形发生率增加。Wilson认为导致出生缺陷率增加的因素：遗传因素占25\%，环境因素占10\%，遗传与环境相互作用及原因不明者占65\%。已知人类致畸因素：辐射（原子武器、放射性碘）；感染（巨细胞病毒、水痘病毒、梅毒螺旋体）；母体代谢失衡（酒精中毒、糖尿病、叶酸缺乏）；药物和环境化学物（氨蝶呤、环磷酰胺、二噁英、氯乙烯等）。
\end{enumerate}

\subsection{致癌物分类}

\begin{table}[htbp]
\centering
\caption{IARC致癌物分类}
\begin{tabularx}{\textwidth}{llX}
\hline
\textbf{类别} & \textbf{分级} & \textbf{说明} \\
\hline
1类 & 对人致癌 & 流行病学证据充分（严格设计、排除混杂、有剂量-反应关系、有独立验证）\\
2A类 & 对人很可能致癌 & 动物实验证据充分，对人体致癌证据有限 \\
2B类 & 对人可能致癌 & 动物实验证据不充分或人体证据不足但动物证据充分 \\
3类 & 尚无法分类 & 可疑对人致癌 \\
4类 & 对人很可能不致癌 & — \\
\hline
\end{tabularx}
\end{table}

\begin{defbox}
\textbf{持久性有机污染物（POPs）}：具有持久性、蓄积性、迁移性和高毒性的有机污染物，可产生急慢性毒性及"三致"效应。

\textbf{环境内分泌干扰物（EDCs）}：对维持机体内环境稳态和调节发育过程的体内天然激素生成、释放、转运、代谢、结合、效应造成严重影响的一类外源性物质。与生殖障碍、出生缺陷、发育异常、代谢紊乱及某些癌症（乳腺癌、睾丸癌、卵巢癌等）的发生发展有关。
\end{defbox}

% ----- 2.5 研究方法 -----
\section{环境与健康关系的研究方法}

\subsection{环境流行病学方法}

\begin{keybox}
\textbf{环境流行病学（Environmental Epidemiology）}：应用传统流行病学方法，结合环境与人群健康关系的特点，从宏观上研究外环境因素与人群健康关系的科学。

\textbf{研究特点：}\imp{多因素、弱效应、非特异性、潜隐期}。

\textbf{基本内容：}
\begin{enumerate}[leftmargin=*]
    \item 研究已知环境因素对人群的健康效应
    \item 探索引起健康异常的环境有害因素
    \item 暴露剂量-反应关系的研究
\end{enumerate}
\end{keybox}

\textbf{（一）暴露测量}

\begin{itemize}[leftmargin=*]
    \item \imp{外暴露剂量}：通过测定人群接触的环境介质中某因素的浓度或含量，根据人体接触特征估计个体暴露水平。减少误差的方法：个体空气采样器、扩大样本量、综合考虑多种暴露途径估计总暴露量。
    \item \imp{内剂量（Internal Dose）}：已吸收至人体内的污染物量。通过测定生物材料（血液、尿液等）中污染物或其代谢产物的含量确定。
    \item \imp{生物有效剂量（Biologically Effective Dose）}：经吸收、代谢活化、转运最终到达靶部位或替代性靶部位的污染物量。
\end{itemize}

\textbf{（二）健康效应测量}

\begin{itemize}[leftmargin=*]
    \item 疾病频率测量：发病率、患病率、死亡率，各种疾病专率
    \item 生理生化功能测量：按手段分（生理、生化、分子生物学、遗传学、影像学）；按器官系统分（呼吸、消化、神经系统等）
\end{itemize}

\textbf{（三）因果关系判断标准}

关联强度、关联稳定性、关联时序性、分布符合性、医学和生物学合理性、剂量-反应关系。需注意\imp{混杂因素（Confounding Factor）}的控制。

\begin{table}[htbp]
\centering
\caption{环境流行病学主要研究方法}
\begin{tabularx}{\textwidth}{lXl}
\hline
\textbf{效应类型} & \textbf{研究方法} & \textbf{特点} \\
\hline
\multirow{3}{*}{急性健康效应} & 时间序列研究 & 分析暴露与效应的时间变化关系 \\
& 病例交叉研究 & 以病例自身为对照 \\
& 固定群组研究 & 追踪固定人群暴露与效应 \\
\hline
\multirow{4}{*}{慢性健康效应} & 横断面研究 & 同一时间点调查暴露与疾病 \\
& 生态学研究 & 以群体为单位分析 \\
& 病例对照研究 & 比较病例与对照暴露史 \\
& 队列研究 & 前瞻性追踪暴露与结局 \\
\hline
\end{tabularx}
\end{table}

\begin{tipbox}
\textbf{宣威女性肺癌高发研究}是环境流行病学研究的经典案例。研究发现室内燃煤空气污染（多环芳烃暴露）是主要原因，改造炉灶后肺癌发病率显著下降。启示：流行病学+干预研究的整合范式。
\end{tipbox}

\subsection{环境毒理学方法}

\begin{table}[htbp]
\centering
\caption{流行病学与毒理学方法的比较}
\begin{tabularx}{\textwidth}{lXX}
\hline
\textbf{比较维度} & \textbf{环境流行病学} & \textbf{环境毒理学} \\
\hline
研究对象 & 人群 & 实验动物/细胞 \\
优势 & 资料来自人群，结果不需外推；可用以评价干预效果；可研究对智力、心理、感觉等的影响 & 可精确控制暴露水平和条件；效应指标不受限制；可应用转基因等先进材料 \\
劣势 & 暴露难精确测量；混杂因素多；弱效应难评价；某些危害（致癌）潜伏期太长 & 种属差异（如反应停和无机砷在动物中为阴性但对人确有毒性）；通常高剂量暴露，易夸大效应 \\
结论强度 & 相关性提示 & 因果推断较强 \\
\hline
\end{tabularx}
\end{table}

\begin{tipbox}
环境流行病学和环境毒理学方法具有\imp{互补性}，开展环境与健康研究时必须将两者结合起来，相互补充、相得益彰。

\textbf{环境生物监测：}传统的环境监测主要采用化学或物理学方法测定介质中污染物的含量，而生物监测则能够迅速反映污染物是否能对生物体特别是体内的遗传物质产生影响。环境理化监测结合生物监测是今后环境监测的趋势。包括：①现场生物监测——对环境动植物或微生物进行细胞遗传学或分子毒理学的直接监测；②环境样品的生物监测——收集空气、水和固体环境样品进行毒理学测试。
\end{tipbox}

\subsection{生物标志物}

\begin{defbox}
\textbf{生物标志物（Biomarker）}：生物体内发生的与发病机制有关联的关键事件的指示物，是机体由于接触各种环境因素引起器官、细胞和亚细胞的生理、生化、免疫和遗传等任何可测定的改变。

\textbf{分类：}
\begin{enumerate}[leftmargin=*]
    \item \imp{暴露生物标志物}——包括内剂量和生物有效剂量标志
    \item \imp{效应生物标志物}——机体内可测定的生理、生化或其他方面的改变
    \item \imp{易感性生物标志物}——能指示机体接触某种特定环境因子时反应能力的标志
\end{enumerate}

\textbf{分子生物标志（Molecular Biomarkers）}：研究外来因子与生物大分子（核酸、蛋白质）相互作用引起的一切分子水平上的改变。以分子生物标志建立的分子流行病学可准确反映暴露与反应的关系，对早期预测损害、评价危险度有重大意义。

\textbf{在环境流行病学中的应用：}暴露的精确测量；揭示早期生物效应；判断宿主易感性。
\end{defbox}

% ----- 2.6 预防对策 -----
\section{环境与健康问题的预防对策}

主要原则：\imp{控制危险因素} → \imp{减少人群接触} → \imp{保护易感人群} → \imp{提升自我防护能力}

主要措施：组织措施、规划措施、技术措施、教育措施。

% ----- 2.7 标准体系 -----
\section{环境与健康标准体系}

我国环境与健康标准体系分为由\imp{环境保护部门}牵头制定的\imp{环境保护标准体系}和由\imp{卫生部门}牵头制定的\imp{环境卫生标准体系}。

\textbf{环境保护标准体系：}以保护人的健康和生存环境、防止生态环境遭受破坏为目的。按级别分三级：国家标准、地方标准（效率高于国家标准）、行业标准。按内容分五类：①环境质量标准（体系核心）；②污染物排放（控制）标准；③环境基础标准；④环境监测分析方法标准；⑤环境样品标准。按效力分：强制性环境标准和推荐性环境标准（不具有法律强制力）。

\textbf{环境卫生标准体系：}以保护人群身体健康为直接目的。包括环境卫生专业基础标准和环境卫生单项标准（室内空气、生活饮用水、公共场所、化妆品等）。

\begin{defbox}
\textbf{基准（Criterion）}：根据环境中有害物质与机体之间的剂量-反应关系，考虑敏感人群和暴露时间而确定的对健康不会产生直接或间接有害影响的相对安全剂量。

\textbf{标准（Standard）}：以保护人群健康为直接目的，对环境中有害影响的因素提出的限量要求以及实现这些要求所规定的相应措施。具有\imp{法律强制性}。
\end{defbox}

\begin{keybox}
\textbf{标准与基准的区别：}

\begin{table}[H]
\centering
\caption{基准与标准的比较}
\begin{tabularx}{\textwidth}{lXX}
\hline
\textbf{比较维度} & \textbf{基准（Criterion）} & \textbf{标准（Standard）} \\
\hline
定义 & 根据剂量-反应关系，考虑敏感人群和暴露时间，确定的对健康不产生有害影响的最大剂量或浓度 & 以保护人群健康为直接目的，对环境中有害因素提出的限量要求及实现这些要求规定的措施 \\
关系 & 标准的科学依据（核心） & 基准内容的实际体现 \\
法律效力 & 不考虑经济、社会、技术因素，无法律效力 & 考虑经济、社会、技术因素，有法律效力 \\
\hline
\end{tabularx}
\end{table}

\begin{itemize}[leftmargin=*]
    \item 基准是纯科学数据，无法律效力
    \item 标准基于基准制定，但需考虑经济、技术可行性
    \item 标准具有法律强制力，是执法依据
    \item 原则上标准限值 $\leq$ 基准值（更严格），但可有例外
\end{itemize}
\end{keybox}

\textbf{制定环境卫生标准的原则：}
\begin{enumerate}[leftmargin=*]
    \item 不引起急性或慢性中毒及潜在的远期危害
    \item 对主观感觉无不良影响
    \item 对人体健康无间接不良影响
    \item 选用最敏感指标（最敏感人群、最敏感观察指标、最灵敏方法）
    \item 经济合理和技术可行
\end{enumerate}

\begin{tipbox}
\textbf{环境卫生标准制定的方法（$\approx$基准的研究方法）：}
\begin{enumerate}[leftmargin=*]
    \item \imp{环境毒理学方法}：明确污染物的作用性质、特点及阈剂量。需考虑安全系数以调整：动物与人种属差异、人群易感性差异、高剂量短时间与低剂量长期暴露差异、小样本外推大群体差异
    \item \imp{感官机能影响的测定}：确定对眼、口腔、上呼吸道的刺激作用阈和异常颜色、气味的阈浓度。在确保安全条件下人体实验比动物实验意义更大
    \item \imp{环境流行病学研究}：验证阈浓度，分析污染与危害的相关程度，评价最高容许浓度的安全程度
    \item 其他：人类受控实验（志愿者）、人类负荷量测定、数学模型、危险度评价
\end{enumerate}
\end{tipbox}

% ----- 2.8 健康风险评价 -----
\section{健康危险度评价}

\begin{keybox}
\textbf{健康危险度评价（Health Risk Assessment, HRA）}：按一定准则，应用毒理学研究和流行病学调查等资料，系统科学地表征有害环境因素暴露对人类和生态的潜在损害作用，并对产生这种损害作用证据的强度或充分性进行评定，对危险性评估相关的不确定性进行评价。

\textbf{特点：}①健康保护观念的转变——安全是相对的，不可能将有害健康的污染物完全清除，只能逐步控制使之处于可接受水平；②把环境污染对健康的影响定量化。

\textbf{四个阶段：}
\begin{enumerate}[leftmargin=*]
    \item \imp{危害鉴定（Hazard Identification）}——定性评价，判断某物质是否对健康有害。依据主要来自流行病学（说服力最强）和毒理学
    \item \imp{暴露评价（Exposure Assessment）}——关键步骤，评估人群暴露的强度、频率和持续时间
    \item \imp{剂量-反应关系评价（Dose-Response Assessment）}——核心，量化剂量与效应关系。有阈化学物采用NOAEL法或基准剂量法推导参考剂量；无阈化学物通过数学模型外推低剂量风险
    \item \imp{风险特征分析（Risk Characterization）}——综合前三阶段，分析判断人群发生某种危害的可能性大小，并阐述不确定性
\end{enumerate}

\textbf{可接受危险度：}社会公认的可接受不良健康效应的概率。有阈化学物参考剂量对应的可接受危险度定为$10^{-6}$。

\textbf{致癌强度系数（CPF）：}实验动物或人终生暴露于剂量为每日每公斤体重1mg致癌物时的终生超额患癌危险度，为剂量-反应曲线斜率的95\%上限。
\end{keybox}

\begin{tipbox}
\textbf{致癌物评价的不确定性来源：}
\begin{enumerate}[leftmargin=*]
    \item 种属差异——不同物种间生理和代谢差异
    \item 高剂量向低剂量外推——高剂量下的反应不一定适用于低剂量
    \item 较短染毒时间向长期持续接触外推
    \item 暴露环境不同——实验暴露条件与人类实际暴露不同
    \item 毒代/毒效动力学资料不足
    \item 数据质量和可靠性问题
\end{enumerate}
\end{tipbox}

% ----- 2.9 复习题 -----
\section{复习题}

\begin{enumerate}[leftmargin=*, itemsep=8pt]
    \item \textbf{【名词解释】}生态系统；食物链；生物放大作用；生物半减期
    \item \textbf{【名词解释】}人群健康效应谱；高危人群；Hormesis效应
    \item \textbf{【名词解释】}生物标志物；环境内分泌干扰物；环境基因交互作用
    \item \textbf{【名词解释】}健康危险度评价（HRA）；可接受危险度
    \item \textbf{【名词解释】}基准与标准；环境生物监测
    \item \textbf{【名词解释】}无阈值化合物与阈值化合物（单阈值与双阈值）
    \item \textbf{【简答题】}简述人群健康效应谱的五级分布规律及其公共卫生学意义。
    \item \textbf{【简答题】}试比较环境流行病学与环境毒理学在研究环境污染健康效应中的优缺点。
    \item \textbf{【简答题】}简述健康危险度评价的基本内容和四个阶段。
    \item \textbf{【简答题】}简述制定环境卫生标准的原则及其制定方法。
    \item \textbf{【简答题】}简述环境多因素联合作用的五种类型及其定义，并说明研究联合作用的意义。
    \item \textbf{【非标准答案题】}试论生物放大作用对环境污染物健康危害评估的启示。
\end{enumerate}

\subsection*{参考答案要点}

\begin{enumerate}[leftmargin=*, itemsep=5pt]
    \item \textbf{生态系统}：在一定空间范围内，由生物群落及其环境组成，借助物质流、能量流、物种流和信息流联结的稳态系统。\textbf{食物链}：不同营养级生物逐级被吞食的链锁关系。\textbf{生物放大作用}：重金属和难降解有机物通过食物链转移到高位营养级生物体内，浓度逐级放大的现象。
    \item \textbf{人群健康效应谱}：人群暴露于污染物时产生的不良反应在性质、程度和范围上表现的征象，从弱到强分为：负荷增加→生理代偿→亚临床状态→临床疾病→死亡。呈金字塔形分布。\textbf{高危人群}：敏感人群和高暴露人群。
    \item \textbf{生物标志物}：可测量的反映暴露、效应或易感性的生物学指标。\textbf{EDCs}：具有类似激素作用，干扰内分泌功能的外源性物质。
    \item HRA：对有害环境因素作用于特定人群的有害健康效应进行综合定性、定量评价的过程，包括危害鉴定、暴露评价、剂量-反应关系评价和风险特征分析四个阶段。
    \item \textbf{基准}：基于剂量-反应关系的纯科学数据，无法律效力。\textbf{标准}：以保护健康为目的的限量要求，具有法律强制力。
    \item \textbf{无阈化合物}：任何剂量均可能有害，无安全剂量（如遗传毒性致癌物）。\textbf{阈值化合物}：仅达到或超过阈剂量才产生效应。
    \item （答题要点：五级分布的具体内容；呈金字塔分布；公共卫生学意义包括早期发现高危人群、指导标准制定、评估干预效果等）
    \item （答题要点：流行病学——优势是结果来自人群无需外推，不足是暴露难精确测量、混杂因素多；毒理学——优势是可精确控制暴露条件，不足是种属差异和高剂量外推问题）
    \item （答题要点：危害鉴定、暴露评价、剂量-反应关系评价、风险特征分析四阶段）
    \item （答题要点：不引起急慢性中毒及远期危害；对主观感觉无不良影响；对健康无间接不良影响；选用最敏感指标；经济合理技术可行）
    \item 四种类型：相加（联合=各单和）、协同（联合>各单和）、增强（无毒物增强有毒物）、拮抗（联合<各单和）。
    \item （开放性问题，从污染物通过食物链传递浓缩、不同营养级暴露风险差异、对风险评价模型的影响等方面展开。）
\end{enumerate}

\newpage

% =====================================================================
%  CHAPTER 3: 室外空气环境与健康 【重点章节】
% =====================================================================
\chapter{室外空气环境与健康  【重点章节】}
\setcounter{chapter}{3}
\chapterintro{
本章为\keyword{重点考察章节}（6学时）。掌握大气层的结构及卫生学意义；掌握大气污染对健康的直接和间接危害（温室效应、臭氧层破坏、酸雨等）；掌握影响大气污染物浓度的因素；掌握主要大气污染物（SO$_2$、颗粒物、NOx、O$_3$等）的健康效应；理解AQI、大气卫生标准等。
}

% ----- 3.1 大气特征 -----
\section{大气的特征及卫生学意义}

\subsection{大气层结构}

\begin{table}[htbp]
\centering
\caption{大气圈五层结构及特征}
\begin{tabularx}{\textwidth}{l>{\raggedright\arraybackslash}p{2.5cm}>{\raggedright\arraybackslash}X}
\hline
\textbf{层次} & \textbf{高度范围} & \textbf{主要特征} \\
\hline
\imp{对流层} & 10--16 km & 气温水平/垂直运动；集中3/4大气质量；气象变化主要发生层；\textbf{与人类健康关系最密切} \\
\imp{平流层} & 16--55 km & 含\keyword{臭氧层}（20--30 km），吸收紫外线；同温层气温恒定 \\
中间层 & 55--85 km & 高度升高，气温下降 \\
热成层 & 85--800 km & 高度升高，气温上升 \\
逸散层 & $>$800 km & 大气稀薄，气体分子逸散 \\
\hline
\end{tabularx}
\end{table}

\begin{defbox}
\textbf{电离层（Ionosphere）}：约60--1000 km，覆盖中间层顶部、热成层和逸散层底部。由于太阳和宇宙射线的作用，大部分空气分子发生电离，具有较高密度的带电粒子，能将电磁波反射回地球，对全球无线电通讯有重大意义。
\end{defbox}

\subsection{大气的组成}

\begin{itemize}[leftmargin=*]
    \item \imp{干洁空气}：不含水汽和气溶胶
    \item \imp{水汽}
    \item \imp{气溶胶}：液体或固体微粒均匀分散在气体中形成的相对稳定的悬浮体系
\end{itemize}

\begin{defbox}
\textbf{雾与霾的区分：}
\begin{itemize}[leftmargin=*]
    \item \textbf{雾}：由大量悬浮在近地面空气中的微小水滴或冰晶组成，影响能见度，是水汽凝结（凝华）的产物。相对湿度 $>$ 90\%
    \item \textbf{霾}：由于空气中悬浮的大量颗粒物所致的水平能见度降低到10 km以下。相对湿度 $<$ 80\%
    \item 介于80\%--90\%之间为雾和霾共同作用的结果
\end{itemize}
\end{defbox}

\subsection{空气离子化}

\begin{defbox}
\textbf{空气离子化（Air Ionization）}：大气分子或原子形成带电荷的正负离子过程。\keyword{负离子}具有清洁空气、改善健康的积极作用。

\textbf{轻离子（Light Ion）}：带电气体分子，包括阳离子和阴离子。

\textbf{重离子（Heavy Ion）}：轻离子与空气中的悬浮颗粒或水滴结合形成。

清洁空气中轻离子浓度高，污染空气中重离子浓度高。空气中重离子数与轻离子数之比 $<$ 50 时，空气较清洁。
\end{defbox}

\subsection{太阳辐射的卫生学意义}
\begin{itemize}[leftmargin=*]
    \item \imp{紫外线（UV, 200--400 nm）}：
    \begin{itemize}
        \item UV-A（320--400 nm）：色素沉着作用
        \item UV-B（290--320 nm）：红斑作用、抗佝偻病作用
        \item UV-C（200--290 nm）：杀菌作用
        \item 过量可致皮肤癌
    \end{itemize}
    \item \imp{可见光（400--760 nm）}：综合作用于高级神经系统，提高视觉和代谢能力，平衡兴奋和镇静作用
    \item \imp{红外线（760--30000 nm）}：热效应；适量可促进新陈代谢和细胞增殖，过强可引起日射病和红外线白内障
\end{itemize}

\textbf{气象因素}（气温、气湿、气压和气流）对机体冷热感觉、体温调节、心血管功能、神经功能、免疫和新陈代谢功能有调节作用。

% ----- 3.2 大气污染 -----
\section{大气污染及污染物的转归}

\subsection{大气污染的来源}

\begin{enumerate}[leftmargin=*]
    \item \imp{工业污染}：燃料燃烧，排放SO$_2$、NOx、颗粒物等
    \item \imp{生活炉灶和采暖锅炉}：燃煤产生SO$_2$、烟尘
    \item \imp{交通运输}：汽车尾气（NOx、CO、VOCs、颗粒物）
    \item \imp{其他}：地面扬尘、意外事件等
\end{enumerate}

\subsection{影响大气污染物浓度的因素}

\begin{keybox}
\textbf{影响大气污染物浓度的三大因素：}

\textbf{一、污染源的排放情况}
\begin{enumerate}[leftmargin=*]
    \item \imp{排放量}：决定大气污染程度的最基本因素
    \item \imp{与污染源的距离}：
    \begin{itemize}
        \item \textbf{烟波着陆点}：烟气自烟囱排出后，向下风侧扩散稀释，接触地面的点。污染物浓度烟波着陆点最大，下风侧浓度随距离增大而降低
    \end{itemize}
    \item \imp{有效排出高度}：烟囱高度 + 烟气抬升高度。排出高度越高，污染物稀释程度越大。下风侧污染物最高浓度与有效排出高度的平方成反比
\end{enumerate}

\textbf{二、气象条件}
\begin{enumerate}[leftmargin=*]
    \item \imp{风和湍流}：水平输送和垂直扩散污染物；湍流强有利于稀释扩散

    \item \imp{风向频率图（风玫瑰图）}：按8或16个方位将一定时期的风向频率绘制成图。风向频率 = 某方向吹来风的次数 / 各方向吹来风的总次数 $\times$ 100\%。\textbf{主导风向}为风向频率最大的风向，其下风向地区污染最严重。只有一个主导风向时，居民区设在主导风向上风向，工业区设在下风向；有多个主导风向时，居民区设在最小风向的下风向，工业区设在其上风向。

    \item \imp{大气稳定度}：\\
    \begin{tabular}{ll}
    $\gamma > \gamma_d$（气温直减率 > 干绝热直减率）& 不稳定状态，波浪型 \\
    $\gamma = \gamma_d$ & 中性状态 \\
    $\gamma < \gamma_d$ & 稳定状态，锥型 \\
    $\gamma < 0$ & \keyword{逆温}，极稳定状态，扇型 \\
    \end{tabular}

    \item \imp{逆温（Temperature Inversion）}：大气温度随高度升高而上升，上层温度高于下层，极不利于污染物扩散。常见类型：
    \begin{itemize}
        \item \textbf{辐射逆温}：无风少云的夜晚，地面通过辐射失去热量而冷却，近地空气随之冷却，上层降温较慢
        \item \textbf{下沉逆温}：空气压缩增温，上层空气下沉落入高气压团中受压变热
        \item \textbf{地形逆温}：局部地理条件形成（如盆地和山谷），冷空气沿山坡聚集在谷底，暖空气受挤压抬升
    \end{itemize}

    \item \imp{气压、气湿、降水}：降水能清除部分污染物；气湿大加重局部污染，利于二次污染物形成（如SO$_2$溶于水易被氧化为硫酸，促进酸雨形成）
\end{enumerate}

\textbf{三、地形条件}
\begin{enumerate}[leftmargin=*]
    \item 山地和谷地：容易形成逆温，不利扩散
    \item \imp{海陆风}：白天升温陆地$>$水面，形成海风；夜晚散热陆地$>$水面，形成陆风。若污染源在岸边，白天可能污染岸上居民区
    \item \imp{城市热岛效应}：人口密集的城市热量散发远大于郊区，城市热空气上升，郊区冷空气补充，把郊区排放的污染物引入城市，造成城市大气污染
\end{enumerate}
\end{keybox}

\subsection{大气污染物的种类}

\textbf{按属性分：}物理性（噪声、电离辐射、电磁辐射等）、化学性和生物性（经空气传播的病原微生物和植物花粉等）三类，其中以化学性污染种类最多、污染范围最广。

\textbf{按存在状态分：}气态和气溶胶。气态污染物包括气体和蒸汽；气溶胶中分散的各种颗粒即大气颗粒物。

\subsection{大气物理性污染与生物性污染}

\begin{defbox}
\textbf{物理性污染：}
\begin{itemize}[leftmargin=*]
    \item \imp{噪声污染}：影响听觉系统（耳鸣、听阈位移、高频听力丧失、耳聋）、神经系统、心血管系统
    \item \imp{光污染}：白亮污染、人工白昼、彩光污染，影响视觉系统、生物钟
    \item \imp{放射性辐射污染}：肺癌（氡）
    \item \imp{电磁辐射污染}
\end{itemize}

\textbf{生物性污染：}
\begin{itemize}[leftmargin=*]
    \item 微生物：呼吸道传染疾病（集中在冬春季）
    \item 变应原：植物花粉、真菌孢子等，引起哮喘、过敏
    \item 其他生物质：柳絮、梧桐叶等生物性尘埃，引起过敏、哮喘、鼻炎、皮肤瘙痒
\end{itemize}
\end{defbox}

\subsection{二次污染}

\begin{defbox}
\textbf{二次污染（Secondary Pollution）}：某些污染物进入大气环境后，本身没有消失，一旦转移后环境发生变化，该污染物可回到原先环境造成二次污染。

\textbf{二次污染物}：由一次污染物在大气中经化学或光化学作用转变成的毒性更大的化学物质。如SO$_2$ → 硫酸雾，NOx + VOCs + 日光 → 光化学烟雾（O$_3$、PAN）。

\textbf{大气污染物的转归：}①自净——扩散和沉降、氧化和中和、植物吸附；②转移——下风侧转移、向其他环境介质转移、向平流层转移；③转化——形成二次污染和二次污染物。
\end{defbox}

% ----- 3.3 大气污染的健康影响 -----
\section{室外空气污染对健康的影响}

\subsection{直接危害}

\begin{keybox}
\textbf{（一）急性危害}

1. \imp{烟雾事件}：

\begin{table}[H]
\centering
\caption{煤烟型烟雾与光化学烟雾的比较}
\begin{tabularx}{\textwidth}{lXX}
\hline
\textbf{比较要点} & \textbf{煤烟型烟雾（London Smog）} & \textbf{光化学烟雾（Photochemical Smog）} \\
\hline
主要污染物 & SO$_2$、颗粒物 & O$_3$、PAN、醛类（光化学氧化剂） \\
污染来源 & 燃煤 & 汽车尾气（NOx + VOCs） \\
气象条件 & 冬季、早晨、低温、辐射逆温、高湿有雾 & 夏秋季、午后、高温、下沉逆温、低湿、紫外线强烈 \\
地形 & 山谷或盆地 & 三面环山，西面临海（如洛杉矶） \\
主要健康效应 & 刺激呼吸道 & 刺激眼睛和呼吸道 \\
\hline
\end{tabularx}
\end{table}

2. \imp{生产事故}：如印度博帕尔异氰酸甲酯泄漏事件（1984）。

\textbf{（二）慢性危害}
\begin{enumerate}[leftmargin=*]
    \item 影响呼吸系统功能：反复刺激→咽炎、喉炎、气管炎→气道狭窄、纤毛损伤、肺泡扩张→COPD；SO$_2$、O$_3$、NOx吸附于颗粒物成为变态反应原，引起支气管哮喘
    \item 心血管疾病：动脉粥样硬化、缺血性心脏病、心律失常、心衰、心搏骤停
    \item 降低免疫功能
    \item 引起变态反应
    \item 致癌作用（如多环芳烃与肺癌）
\end{enumerate}
\end{keybox}

\subsection{间接危害}

\begin{keybox}
\textbf{1. 温室效应（Greenhouse Effect）}
\begin{itemize}[leftmargin=*]
    \item 原因：CO$_2$、CH$_4$等温室气体吸收地表热辐射，使大气增温
    \item 危害：海平面上升；促进酸雨形成；热相关疾病增加；虫媒传染病流行范围扩大
\end{itemize}

\textbf{2. 臭氧层破坏}
\begin{itemize}[leftmargin=*]
    \item 原因：CFCs（氯氟烃）、N$_2$O、CCl$_4$、哈龙等消耗臭氧层物质
    \item 危害：皮肤癌、白内障发生率增加；影响生态系统
\end{itemize}

\textbf{3. 酸雨}
\begin{itemize}[leftmargin=*]
    \item 定义：降水pH值 $<$ 5.6即为酸雨
    \item 我国以\keyword{硫酸型酸雨}为主，主要因燃煤排放SO$_2$和NOx
    \item 危害：土壤和植物损害；影响水生生态系统；腐蚀建筑物和文物；影响人类健康
    \item 防治：控制SO$_2$和NOx排放
\end{itemize}

\textbf{4. 大气棕色云团（ABC）}
\begin{itemize}[leftmargin=*]
    \item 区域性大气污染现象，由颗粒物和气体污染物共同形成
\end{itemize}
\end{keybox}

\subsection{主要大气污染物及其健康影响}

\begin{table}[htbp]
\centering
\caption{主要大气污染物及其健康效应}
\begin{tabularx}{\textwidth}{lX}
\hline
\textbf{污染物} & \textbf{主要健康效应} \\
\hline
\imp{颗粒物（PM）} & 我国多数城市的\keyword{首要污染物}。呼吸系统：堵塞、刺激肺泡壁、氧化应激炎症反应；心血管：干扰自主神经→心律不齐、诱发血栓、刺激呼吸道释放细胞因子损伤血管。粒径$>$5$\mu$m沉积上呼吸道，$<$5$\mu$m进入细支气管和肺泡 \\
\imp{SO$_2$} & 阻滞于上呼吸道→亚硫酸盐/硫酸氢盐；气管和支气管收缩，气道阻力↑；吸附SO$_2$的颗粒物是变态反应原→支气管哮喘；促癌（增加BaP致癌性） \\
\imp{NOx} & NO$_2$难溶于水，主要作用于深部呼吸道、细支气管和肺泡；形成亚硝酸盐和硝酸盐→高铁血红蛋白→组织缺氧；与SO$_2$、O$_3$有相加和协同作用 \\
\imp{O$_3$} & 二次污染物（占光化学氧化剂90\%以上），光化学烟雾的指示物；体内→超氧阴离子→呼吸道上皮脂质过氧化→炎症；激活NF-$\kappa$B \\
\imp{CO} & 与Hb亲和力比O$_2$大200--300倍→HbCO→降低血液携氧能力→阻碍氧释放→组织缺氧；还可与肌红蛋白、细胞色素氧化酶及P450结合；急性中毒以神经系统症状为主；\textbf{胎儿对CO毒性比成人敏感}——孕妇吸烟可致低体重儿、围产期死亡增高及婴幼儿神经行为障碍 \\
\imp{Pb（铅）} & 全身性毒物，90\%贮存于骨骼。神经系统（神经衰弱→周围神经炎）；造血系统（抑制Hb合成→贫血）；可透过胎盘屏障影响胚胎发育。\textbf{儿童为敏感人群}（户外活动多、摄入量大、吸收率高、血脑屏障发育不完全）——铅选择性蓄积于脑海马部位，损害神经细胞形态和功能，表现为注意力不集中、记忆力降低、学习能力下降等。交警为铅职业接触高危人群 \\
\imp{PAHs（多环芳烃）} & 致癌（如苯并[a]芘与肺癌） \\
\imp{二噁英} & 主要来源于垃圾焚烧、工业生产；吸附于颗粒物→沉降到水体和土壤→通过食物链富集进入人体。食物是人体内二噁英的主要来源 \\
\hline
\end{tabularx}
\end{table}

\begin{defbox}
\textbf{空气动力学等效直径（Dp）}：在气流中，被测颗粒的空气动力学效应与密度为1的球形标准颗粒物相同时，该标准颗粒物粒径即被测颗粒物的Dp。可表达颗粒物在空气中的停留时间、沉降速度、进入呼吸道的可能性及在呼吸道的沉降部位。

\textbf{TSP（总悬浮颗粒物）}：粒径 $\leq$ 100 $\mu$m的颗粒物，包括液体、固体或液体与固体结合并悬浮在空气介质中的颗粒。

\textbf{IP（可吸入颗粒物）/ PM$_{10}$}：粒径 $\leq$ 10 $\mu$m，能进入人体呼吸道，长期漂浮在空气中。

\textbf{PM$_{2.5}$（细颗粒物）}：粒径 $\leq$ 2.5 $\mu$m，可进入肺泡和血液循环，易吸附有毒有机物和重金属元素，对健康危害极大。

\textbf{PM$_{0.1}$（超细颗粒物）}：粒径 $\leq$ 0.1 $\mu$m。
\end{defbox}

% ----- 3.4 调查与监测 -----
\section{大气污染和气候变化对健康影响的调查与监测}

\subsection{环境空气质量功能分区与标准分级}

\begin{keybox}
\textbf{环境空气质量功能分区：}
\begin{itemize}[leftmargin=*]
    \item \imp{一类区}：自然保护区、风景名胜区和其他需要特殊保护的地区 → 执行\imp{一级标准}（长期接触不发生任何危害影响）
    \item \imp{二类区}：居住区、商业交通居民混合区、文化区、工业区和农村地区 → 执行\imp{二级标准}（长期和短期接触不发生伤害）
\end{itemize}

\textbf{污染物浓度限值类型：}
\begin{itemize}[leftmargin=*]
    \item \imp{1小时平均浓度}：任何1h内平均浓度的最高容许值。针对能在短期内产生刺激、过敏或中毒等急性危害的物质
    \item \imp{24小时平均浓度（日平均）}：任何自然日24h平均浓度的最高容许值
    \item \imp{年平均浓度}：任何日历年内各日平均浓度的算术均值的最高容许值。针对有慢性作用的物质，防止慢性和潜在危害
\end{itemize}
\end{keybox}

调查内容：①污染源调查；②污染状况监测；③人群健康调查。

\subsection{采样布点方法}

\begin{table}[htbp]
\centering
\caption{大气污染监测采样布点方法}
\begin{tabularx}{\textwidth}{lX}
\hline
\textbf{污染源类型} & \textbf{布点方法} \\
\hline
\imp{点源污染} & ①四周布点：以污染源为中心，8个方位不同距离同心圆布点；②扇形布点：主导方向下风侧划3--5个方位布置；③捕捉烟波布点：随烟波方向变动布点（不固定） \\
\imp{面源污染} & ①按城市功能分区布点（每种类型2--3个采样点+清洁对照点）；②几何状布点（方形或三角形网格）；③综合因素布点 \\
\imp{线源污染} & 采样口离地面高度2--5 m，距道路边缘$\leq$20 m \\
\hline
\end{tabularx}
\end{table}

% ----- 3.5 AQI -----
\section{大气质量标准}

\begin{defbox}
\textbf{日平均最高容许浓度}：任何一天内多次测定的平均浓度的最高容许值。

\textbf{一次最高容许浓度}：任何一次短时间采样测量结果的最高容许值。
\end{defbox}

\begin{tipbox}
\textbf{AQI（空气质量指数）}：属于分段线性函数型大气质量指数。参数包括SO$_2$、NO$_2$、PM$_{10}$、PM$_{2.5}$、CO、O$_3$共6项。

\textbf{空气质量分指数（IAQI）}：各污染物的分指数。AQI取各IAQI中的最大值，该污染物即为\imp{首要污染物}。

\textbf{确定方法：}
\begin{itemize}[leftmargin=*]
    \item AQI $>$ 50时，IAQI最大的污染物为首要污染物（若多项并列则并列为首要污染物）
    \item IAQI $>$ 100的污染物为\imp{超标污染物}
\end{itemize}

\textbf{优缺点：}优点——突出当日单个污染物的作用，便于发现首要污染物及其对健康的影响；缺点——对各污染物指数无综合，仅给出最大分指数，不能反映其他污染物的综合影响。
\end{tipbox}

\section{气候变化与极端天气的健康影响}

\subsection{气候变化（Climate Change）}

\begin{keybox}
\textbf{气候变化对健康的直接影响：}
\begin{itemize}[leftmargin=*]
    \item \imp{高温热浪}：可改变人体生理功能，增加中暑、热射病、急性肾衰竭等急性病风险
    \item \imp{洪涝与风暴潮}：造成溺水和创伤，导致死亡率和伤残率增加
    \item \imp{野火燃烧和沙尘天气}：通过携带大气颗粒物等污染物直接影响健康
\end{itemize}

\textbf{气候变化对健康的间接影响：}
\begin{enumerate}[leftmargin=*]
    \item \imp{传染病风险加剧}：温度升高与降水模式改变使昆虫传播媒介和病原体的地理分布范围扩大、繁殖季节延长，导致病毒传播到纬度与海拔更高的地区
    \item \imp{粮食安全}：干旱、洪涝、高温热浪导致农作物产量减少及粮食营养成分下降
    \item \imp{水资源短缺}：影响饮用水供应与水质，导致卫生条件恶化，引发肠道传染病
    \item \imp{过敏性疾患增加}：气候变暖使真菌孢子、花粉等浓度增高
\end{enumerate}
\end{keybox}

\subsection{极端天气}

\begin{defbox}
\textbf{极端天气事件}：显著偏离正常气候状况，对人类社会和自然生态系统构成严重威胁的高强度、低频次天气现象。特征为强度高、影响范围广、持续时间长。

\begin{itemize}[leftmargin=*]
    \item \imp{高温热浪}：增加中暑、热射病、热痉挛等热相关疾病风险
    \item \imp{寒潮}：增加心血管疾病发作风险，削弱皮肤和呼吸道黏膜免疫屏障
    \item \imp{极端强降水}：引发山洪、内涝，造成人员伤亡；污染清洁水源，导致病原微生物和媒介生物快速孳生，造成传染病暴发流行
    \item \imp{洪涝}：引发心理疾病（抑郁、焦虑）；水源污染和细菌孳生（霍乱、痢疾等）引发肠道和呼吸道感染；破坏公共设施造成救护延误，增加慢性非传染性疾病死亡风险
    \item \imp{干旱}：粮食减产，受灾人群营养不良、营养缺乏风险增加
    \item \imp{野火}：烟雾中有害颗粒物和气体引发或加剧哮喘、COPD等呼吸系统疾病，增加心血管疾病风险
    \item \imp{沙尘暴}：加重呼吸系统疾病，引起眼部刺激和结膜炎，增加过敏反应
    \item \imp{台风}：机体损伤（软组织挫伤、颅脑损伤、骨折），增加介水传染病和食源性传染病风险
\end{itemize}
\end{defbox}

\section{大气卫生防护措施}

\begin{itemize}[leftmargin=*]
    \item \textbf{规划措施}：合理安排工业布局和城镇功能分区；加强绿化；加强局部污染源管理
    \item \textbf{工艺措施}：控制燃煤污染（清洁能源、集中供热）；加强废气治理（除尘、脱硫、脱硝）
\end{itemize}

% ----- 3.7 复习题 -----
\section{复习题}

\begin{enumerate}[leftmargin=*, itemsep=8pt]
    \item \textbf{【名词解释】}光化学烟雾；逆温；酸雨
    \item \textbf{【名词解释】}TSP；PM$_{10}$；PM$_{2.5}$
    \item \textbf{【名词解释】}一次污染物与二次污染物（大气中）
    \item \textbf{【名词解释】}温室效应
    \item \textbf{【简答题】}试述大气稳定度对大气污染扩散的影响。
    \item \textbf{【简答题】}请简述影响大气中污染物浓度的因素。
    \item \textbf{【简答题】}大气污染对健康的直接危害有哪些？间接危害有哪些？
    \item \textbf{【简答题】}臭氧层破坏的原因及其健康危害是什么？
    \item \textbf{【简答题】}简述光化学烟雾与煤烟型烟雾的主要区别。
    \item \textbf{【简答题】}简述酸雨的形成原因、危害及防治对策。
    \item \textbf{【简答题】}简述大气物理性污染和生物性污染的主要类型及健康影响。
    \item \textbf{【简答题】}简述气候变化对健康的直接和间接影响。
    \item \textbf{【简答题】}简述极端天气事件的类型及其健康危害。
    \item \textbf{【简答题】}简述大气卫生防护的主要措施。
    \item \textbf{【非标准答案题】}结合我国当前大气污染现状，谈谈你对大气污染防治与人群健康保护的思考。
\end{enumerate}

\subsection*{参考答案要点}

\begin{enumerate}[leftmargin=*, itemsep=5pt]
    \item \textbf{光化学烟雾}：汽车尾气中NOx和VOCs在日光紫外线照射下经光化学反应生成的浅蓝色烟雾，主要含O$_3$、PAN等强氧化性物质。\textbf{逆温}：大气温度随高度升高而上升，上层气温高于下层，极不利于污染物扩散。\textbf{酸雨}：降水pH值$<$5.6，主要由SO$_2$和NOx转化形成。

    \item \textbf{TSP}：总悬浮颗粒物，粒径$\leq$100$\mu$m。\textbf{PM$_{10}$}：可吸入颗粒物，粒径$\leq$10$\mu$m，能进入呼吸道。\textbf{PM$_{2.5}$}：细颗粒物，粒径$\leq$2.5$\mu$m，可进入肺泡和血液循环。

    \item （参照第1章定义）

    \item 大气中CO$_2$、CH$_4$等温室气体吸收地表热辐射，使大气增温的作用。

    \item （答题要点：$\gamma$与$\gamma_d$的比较→不稳定/中性/稳定状态；逆温条件下的极稳定状态；各种稳定度对应的烟型及其对污染物扩散的影响）

    \item （答题要点：三大因素——污染源排放情况（排放量、高度、距离）、气象条件（风、湍流、大气稳定度、逆温、降水）、地形条件（山地谷地、海陆风、热岛效应））

    \item 直接危害：急性（烟雾事件、生产事故）和慢性（呼吸系统疾病、心血管疾病、免疫功能降低、致癌）。间接危害：温室效应、臭氧层破坏、酸雨、大气棕色云团。

    \item 原因：CFCs、N$_2$O、CCl$_4$等消耗臭氧层物质释放。危害：皮肤癌和白内障增多，影响生态系统。

    \item 光化学烟雾：汽车尾气+日光，含O$_3$、PAN，夏秋季午后多，刺激眼睛。煤烟型：燃煤+逆温，含SO$_2$、烟尘，冬季清晨多，刺激呼吸道。

    \item 成因：燃煤和石油燃烧排放SO$_2$和NOx，在大气中氧化转化为硫酸和硝酸。危害：土壤酸化、水生生态破坏、建筑物腐蚀、健康影响。防治：控制SO$_2$和NOx排放。

    \item 规划措施（合理布局、绿化、局部源管理）和工艺措施（控制燃煤、加强废气治理）。

    \item （开放性问题，可围绕清洁能源、产业结构调整、多部门协作、公众参与等方面展开。）
\end{enumerate}

\newpage

% =====================================================================
%  CHAPTER 4: 室内环境与健康
% =====================================================================
\chapter{室内环境与健康}
\setcounter{chapter}{4}
\chapterintro{
掌握住宅的卫生学意义和基本卫生要求；掌握室内空气污染的来源和特点；掌握室内空气清洁度的评价指标；了解办公场所的卫生学特点及管理要求。
}

% ----- 4.1 住宅卫生 -----
\section{住宅卫生}

\subsection{住宅的卫生学意义}

\begin{defbox}
\textbf{住宅（Residential Building）}：人们为了充分利用自然环境和人为环境因素中的有利作用和防止其不良影响而创造的生活居住环境。
\end{defbox}

\textbf{住宅的卫生学意义：}
\begin{enumerate}[leftmargin=*]
    \item 住宅是人们生活、学习、工作的最重要场所
    \item 住宅的卫生条件和人类健康息息相关
    \item 不良住宅环境对健康不利；良好住宅环境有利于健康
    \item 住宅卫生状况可影响数代人和众多家庭的健康
    \item 住宅环境对健康的影响具有\keyword{长期性和复杂性}
\end{enumerate}

\subsection{住宅的基本卫生要求}

\begin{enumerate}[leftmargin=*]
    \item \imp{小气候适宜}——温度、湿度、气流、热辐射适宜
    \item \imp{采光照明良好}——充足的自然采光和人工照明
    \item \imp{空气清洁卫生}——通风良好，无有害物质积累
    \item \imp{隔音性能良好}——防止噪声干扰
    \item \imp{卫生设施齐全}——上下水、卫生厕所等
    \item \imp{环境安静整洁}——有利于身心健康
\end{enumerate}

\subsection{居室卫生规模}

\textbf{居室卫生规模}是指根据卫生要求提出的居室面积、容积、净高和进深等应有的规模。

\begin{enumerate}[leftmargin=*]
    \item \imp{居室净高}：室内地板到天花板的高度。面积相同时净高越大容积越大，利于采光通风；过低使人压抑且不利于通风散热
    \item \imp{居室面积（居住面积）}
    \item \imp{居室容积}：每个居住者占有的空间容积 = 净高 $\times$ 面积。室内空气中CO$_2$含量是评价空气清洁度的重要指标
    \item \imp{居室进深}：开设窗户的外墙内表面至对面墙内表面的距离。进深与宽度之比不宜 $>$ 2:1，以3:2为宜。\textbf{室深系数}：进深与地板至窗上缘高度之比，一侧采光 $\leq$ 2--2.5，两侧采光 $\leq$ 4--5
\end{enumerate}

\subsection{住宅朝向与间距}

\begin{itemize}[leftmargin=*]
    \item \imp{住宅朝向}：应使主要房间朝南，辅助房间朝北，满足冬季尽量多日照、夏季避免过多日照并利于自然通风
    \item \imp{住宅间距}：前后排建筑物间应保持的最小间隔距离。\textbf{日照间距}：前后两排南向房屋间，为保证后排房屋在冬至日底层获得$\geq$1h满窗日照保持的最小间隔距离
\end{itemize}

\subsection{隔声}

用隔声材料和隔声结构（实体墙板、密封门窗等）阻挡声能传播，将居室相对封闭起来以减少噪声污染。<｜end▁of▁thinking｜>

\subsection{住宅发展新方向}
\imp{健康住宅}：突出健康要素，满足居住者生理、心理和社会多层次需求。\imp{绿色生态住宅}：消耗最少资源和能源、产生最少废弃物，强调节能、节水、节地和治理污染。

\begin{table}[htbp]
\centering
\caption{居室采光评价指标}
\begin{tabularx}{\textwidth}{lXl}
\hline
\textbf{指标} & \textbf{定义} & \textbf{卫生要求} \\
\hline
采光系数 & 室内工作水平面散射光照度与室外全天空散射光照度的百分比 & — \\
投射角 & 室内工作点与采光口上缘连线和水平线夹角 & $\geq$ 27° \\
开角 & 工作点与对侧遮光物上端连线和工作点与采光口上缘连线的夹角 & $\geq$ 4° \\
窗地面积比值 & 采光口窗玻璃面积与室内地面面积之比 & — \\
\hline
\end{tabularx}
\end{table}

\subsection{室内小气候}

\begin{defbox}
\textbf{小气候（Microclimate）}：生活环境中空气的温度、湿度、气流和热辐射等因素对机体热平衡产生明显影响的综合环境条件。
\end{defbox}

\begin{defbox}
\textbf{有效温度（Effective Temperature, ET）}：人体在不同温度、湿度和风速综合作用下产生的热感觉指标。以风速0 m/s、相对湿度100\%、气温17.7°C时产生的温热感为评价标准。

\textbf{校正有效温度（CET）}：将干球温度替换为黑球温度得出的有效温度。

\textbf{湿球-黑球温度指数（WBGT）}：综合反映气温、气湿、气流和热辐射四种物理因素对机体的作用。

\textbf{风冷指数（WCI）}：综合反映寒冷气候下气温和风速对人体温热感的影响，即在低温环境中风速增加产生的冷效应相当于气温下降的度数。

\textbf{热强度指数（TEI）}：根据机体与环境间热交换情况制订的评价指标。
\end{defbox}

\textbf{小气候对人体影响的常用生理指标：}
\begin{itemize}[leftmargin=*]
    \item \imp{皮肤温度（皮温）}：最常用，与温热感觉、脉搏变化基本平行。\textbf{加权平均皮肤温度（WMST）}：测定3--8个点的皮温
    \item \imp{体温}：判断机体热平衡是否受到破坏最直接的指标
    \item 脉搏、出汗量、温热感、热平衡测定
\end{itemize}

% ----- 4.2 室内空气污染 -----
\section{室内空气污染}

\subsection{室内空气污染的来源}

\begin{keybox}
\textbf{室外来源：}
\begin{enumerate}[leftmargin=*]
    \item 室外空气（大气污染物渗入）
    \item 住宅建筑物材料（氡、甲醛等）
    \item 人为带入室内
    \item 相邻住宅污染
    \item 生活用水污染
\end{enumerate}

\textbf{室内来源：}
\begin{enumerate}[leftmargin=*]
    \item \imp{室内燃烧或加热}——燃煤、燃气、烹调油烟
    \item \imp{室内活动}——吸烟、人体代谢产物
    \item \imp{室内建筑装修材料}——甲醛、VOCs、氡
    \item \imp{室内生物性污染}——霉菌、尘螨、细菌
    \item \imp{家用电器}——电磁辐射、臭氧
\end{enumerate}
\end{keybox}

\subsection{室内空气污染的特点}

\begin{enumerate}[leftmargin=*]
    \item 室外污染物对室内空气的污染
    \item 室内外共存同类污染物的叠加暴露
    \item \imp{吸烟}对室内空气污染（最重要的室内污染源之一）
    \item \imp{建筑材料}对室内空气污染（甲醛、氡等）
    \item 空调引起的室内空气污染（"Sick Building"）
\end{enumerate}

\subsection{居室空气清洁度的评价指标}

\begin{table}[htbp]
\centering
\caption{居室空气清洁度评价指标}
\begin{tabularx}{\textwidth}{lX}
\hline
\textbf{指标} & \textbf{意义} \\
\hline
CO$_2$ & 反映室内人员密度和通风状况，$<$0.07\%（清洁） \\
微生物和悬浮颗粒 & 反映室内卫生状况和通风效果 \\
CO & 反映燃烧不完全程度 \\
SO$_2$ & 反映燃煤污染 \\
空气离子 & 负离子浓度反映空气新鲜程度 \\
甲醛 & 反映装修材料污染 \\
\hline
\end{tabularx}
\end{table}

\subsection{保持居室空气清洁度的卫生措施}

\begin{enumerate}[leftmargin=*]
    \item 住宅地段选择（远离污染源）
    \item 建筑材料和装饰材料选择（绿色环保材料）
    \item 合理的住宅平面配置
    \item 合理的住宅卫生规模
    \item 采用改善空气质量的措施（通风、净化）
    \item 改进个人卫生习惯
    \item 合理使用和保养各种设施
    \item 加强卫生宣传教育
\end{enumerate}

% ----- 4.3 病态建筑物综合征 -----
\section{病态建筑物综合征}

\begin{defbox}
\textbf{病态建筑物综合征（Sick Building Syndrome, SBS）}：现代住宅中多种环境因素（如物理、化学因素）联合作用对健康的影响所引起的一种综合征。主要表现为眼、鼻、喉刺激症状，头痛，疲劳，注意力不集中等。与空调系统使用、通风不良、装修材料污染等有关，确切原因尚不清楚。

\textbf{建筑物相关疾病（Building Related Illness, BRI）}：由于人体暴露于建筑物内有害因素（如细菌、真菌、尘螨、氡、一氧化碳、甲醛等）引起的疾病。与SBS的区别在于BRI有明确的病因和典型的临床症状。

\textbf{化学物质过敏症（Multiple Chemical Sensitivity, MCS）}：由于多种化学物质作用于人体多种器官系统，引起多种症状的疾病。特点：复发性、症状呈慢性过程、由低浓度化学污染物质引发、很难找到具体单一对应致病原。
\end{defbox}

% ----- 4.4 办公场所卫生 -----
\section{办公场所卫生}

\begin{defbox}
\textbf{办公场所（Workplace）}：管理或专业技术人员处理某种特定事务的室内工作环境。
\end{defbox}

\textbf{卫生学特点：}
\begin{enumerate}[leftmargin=*]
    \item 办公人员相对集中，流动性小
    \item 办公人员滞留时间长，活动范围小
    \item 分布范围广泛，基本条件和卫生状况相差较大
    \item 存在诸多影响人体健康的不利因素
\end{enumerate}

\textbf{基本卫生要求：}用地选择合理、采光照明良好、小气候适宜、空气质量良好、环境宽松。

% ----- 4.5 复习题 -----
\section{复习题}

\begin{enumerate}[leftmargin=*, itemsep=8pt]
    \item \textbf{【名词解释】}小气候；有效温度（ET）；病态建筑物综合征（SBS）
    \item \textbf{【名词解释】}采光系数；窗地面积比值
    \item \textbf{【简答题】}简述住宅的基本卫生要求。
    \item \textbf{【简答题】}简述室内空气污染的来源及特点。
    \item \textbf{【简答题】}居室空气清洁度的评价指标有哪些？
    \item \textbf{【简答题】}简述保持居室空气清洁度的卫生措施。
    \item \textbf{【简答题】}简述办公场所的卫生学特点。
\end{enumerate}

\subsection*{参考答案要点}

\begin{enumerate}[leftmargin=*, itemsep=5pt]
    \item \textbf{小气候}：生活环境中空气温度、湿度、气流和热辐射等因素的综合环境条件。\textbf{ET}：人体在不同温度、湿度和风速综合作用下的热感觉指标。\textbf{SBS}：现代住宅中多种环境因素综合作用引起的综合征，表现为眼鼻喉刺激、头痛、疲劳等。

    \item \textbf{采光系数}：室内工作水平面散射光照度与同时室外全天空散射光照度的百分比。\textbf{窗地面积比值}：采光口窗玻璃面积与室内地面面积之比。

    \item 六项基本要求：小气候适宜、采光照明良好、空气清洁卫生、隔音性能良好、卫生设施齐全、环境安静整洁。

    \item 来源分室外（大气渗入、建筑材料、人为带入等）和室内（燃烧加热、室内活动、装修材料、生物性污染、家电）。特点：室外污染渗入、叠加暴露、吸烟污染、建材污染、空调污染。

    \item 主要指标：CO$_2$、微生物和悬浮颗粒、CO、SO$_2$、空气离子、甲醛等。

    \item 八项措施：地段选择、材料选择、平面配置、卫生规模、通风净化、个人习惯、设施保养、卫生宣教。

    \item 四项特点：人员集中流动性小、滞留时间长活动范围小、分布广条件差异大、存在多种不利因素。
\end{enumerate}

\newpage

% =====================================================================
%  CHAPTER 5-10 will be appended next...
% =====================================================================

% =====================================================================
%  CHAPTER 5: 水体卫生 【重点章节】
% =====================================================================
\chapter{水体卫生  【重点章节】}
\setcounter{chapter}{5}
\chapterintro{
本章为\keyword{重点考察章节}（6学时）。掌握水质的评价指标（物理、化学、微生物学性状指标）；掌握水体污染来源及各种水体的污染特点；掌握水体自净机制及氧垂曲线；掌握水体富营养化的概念及危害；掌握水体污染的危害；了解水环境标准和水体卫生防护措施。
}

\section{水资源的种类及卫生学特征}

\begin{defbox}
\textbf{水资源（Water Resources）}：全球水量中对人类生存、发展可用的水量，主要指逐年可得到更新的那部分淡水量。
\end{defbox}

\begin{table}[htbp]
\centering
\caption{天然水资源分类及卫生学特征}
\begin{tabularx}{\textwidth}{lXp{3.5cm}}
\hline
\textbf{类型} & \textbf{卫生学特征} & \textbf{主要问题} \\
\hline
\imp{降水} & 水质较好，矿物质含量低；地区分布极不平衡 & 水量无保证；受大气质量影响（如酸雨）；沿海地区降水含碘较多 \\
\imp{地表水} & 溶解土壤矿物质；水量充足；可分为封闭型（湖水、水库水等）和开放型（江河水） & 人为污染是影响水质最主要因素；有丰水期和枯水期变化 \\
\imp{地下水} & 水质清洁，矿物质含量较高；水量较稳定 & 一旦污染，自净极其困难；硬度较高 \\
\hline
\end{tabularx}
\end{table}

\begin{defbox}
\textbf{透水层}：由颗粒较大的砂、砾石组成，能渗水与存水。

\textbf{不透水层}：由颗粒细小致密的黏土层和岩石层构成。

\textbf{浅层地下水}：位于第一个不透水层之上。水质物理性状较好，细菌数较地表水少，但硬度增高，溶解氧减少。我国广大农村使用。

\textbf{深层地下水}：位于第一个不透水层之下。水质透明无色，水温恒定，水量稳定，细菌数很少，但盐类含量高、硬度大。城镇或企业集中式供水常用。

\textbf{泉水}：地下水通过地表缝隙自行涌出。分潜水泉（浅层地下水涌出）和自流泉（深层地下水由裂隙中涌出）。
\end{defbox}

\section{水质的性状和评价指标}

\subsection{物理性状指标}

\begin{table}[htbp]
\centering
\caption{水质物理性状指标}
\begin{tabularx}{\textwidth}{lX}
\hline
\textbf{指标} & \textbf{含义及意义} \\
\hline
水温 & 影响水中生物和化学反应速率 \\
色 & 纯净水无色；有色可能提示污染 \\
臭和味 & 异常臭和味提示有机物分解或化学污染 \\
\imp{浑浊度} & 表示水中悬浮物和胶体物对光线透过时的阻碍程度。1度=1L水中含1mg标准硅藻土 \\
放射性 & 天然和人为放射性物质含量 \\
\hline
\end{tabularx}
\end{table}

\subsection{化学性状指标}

\begin{keybox}
\textbf{主要化学性状指标：}

\begin{enumerate}[leftmargin=*]
    \item \imp{pH值}：天然水pH一般在7.2--8.5之间
    \item \imp{总固体（Total Solid）}：水样蒸发至干后的残留物总量 = 溶解性固体 + 悬浮性固体
    \item \imp{硬度（Hardness）}：溶于水中钙、镁盐类的总含量，以CaCO$_3$（mg/L）表示
    \begin{itemize}
        \item 暂时硬度：水煮沸后能去除的硬度（碳酸盐硬度）
        \item 永久硬度：煮沸后不能去除的硬度（非碳酸盐硬度）
    \end{itemize}
    \item \imp{含氮化合物}——反映有机污染及其自净进程：
    \begin{itemize}
        \item \textbf{有机氮和蛋白氮}：有机氮是有机含氮化合物的总称，蛋白氮是已分解成较简单的有机氮。两者显著增高说明\keyword{水体新近受明显有机性污染}
        \item \textbf{氨氮}：含氮有机物在有氧条件下经微生物分解形成的最初产物，表示\keyword{新近可能有人畜粪便污染}
        \item \textbf{亚硝酸盐氮}：氨在有氧条件下经亚硝酸菌作用形成，是硝化过程的中间产物，含量高说明\keyword{有机物无机化过程尚未完成，污染仍存在}
        \item \textbf{硝酸盐氮}：含氮有机物氧化分解的最终产物，含量高而氨氮、亚硝酸盐氮不高，说明\keyword{过去曾受有机污染，现已完成自净}
    \end{itemize}
    \begin{table}[H]
    \centering
    \caption{三氮组合判断水体污染状况}
    \begin{tabular}{ccc}
    \hline
    \textbf{NH$_3$-N} & \textbf{NO$_2^-$-N / NO$_3^-$-N} & \textbf{污染状况} \\
    \hline
    高 & 低 & \textcolor{keyred}{新近污染} \\
    低 & 高 & \textcolor{darkgreen}{陈旧污染（已自净）} \\
    高 & 高 & 持续污染（过去和新近均有污染） \\
    均高 & — & 过去受污染，目前自净正在进行 \\
    \hline
    \end{tabular}
    \end{table}
    \item \imp{溶解氧（DO）}：溶解在水中的氧含量，清洁水DO较高
    \item \imp{化学耗氧量（COD）}：用强氧化剂氧化水中有机物所消耗的氧量
    \item \imp{生化需氧量（BOD）}：水中有机物在需氧微生物分解时消耗的溶解氧量。\textbf{BOD$_5^{20}$}：20°C培养5日后1L水中减少的溶解氧量。清洁水BOD$_5$一般 $<$ 1 mg/L
    \item \imp{总有机碳（TOC）}：水中全部有机物的含碳量，是评价水体有机物污染程度的综合性指标之一，但不能说明有机污染物的性质
    \item \imp{总需氧量（TOD）}：1L水中还原物质（有机物和无机物）在一定条件下被氧化时所消耗氧的毫升数，数值愈大污染愈严重。TOC和TOD的检测可实现快速自动化
    \item \imp{氯化物}：含量突然增高时表明水可能受到人畜粪便、生活污水或工业废水的污染
    \item \imp{硫酸盐}：含量突然增加表明水可能受到生活污水、工业废水或硫酸铵化肥等污染
\end{enumerate}
\end{keybox}

\subsection{微生物学性状指标}

\begin{defbox}
\textbf{指示菌（Indicator Bacteria）}：具有代表微生物污染总体状况的菌种。
\end{defbox}

\begin{enumerate}[leftmargin=*]
    \item \imp{细菌总数（Bacteria Count）}：1 mL水在普通琼脂培养基中37°C培养24小时后生长的细菌菌落数。
    \item \imp{总大肠菌群（Total Coliforms）}：需氧及兼性厌氧、37°C生长时使乳糖发酵、24小时内产酸产气的革兰氏阴性无芽孢杆菌。是\keyword{粪便污染指示菌}。
    \item \imp{粪大肠菌群（Fecal Coliforms）}：44.5°C培养仍能生长的大肠菌群，更特异地指示粪便污染。
\end{enumerate}

\section{水体的污染来源及特征}

\subsection{污染来源}

\begin{enumerate}[leftmargin=*]
    \item \imp{工业废水}：排放量大、成分复杂、毒性强、难处理
    \item \imp{生活污水}：含有机物、病原微生物、氮磷等
    \item \imp{农业污水}：含农药、化肥、畜禽养殖废弃物
\end{enumerate}

\subsection{各种水体的污染特点}

\begin{keybox}
\begin{enumerate}[leftmargin=*]
    \item \imp{河流}：污染程度取决于\keyword{径污比}（径流量与排入污水量的比值）。
    \item \imp{湖泊、水库}：稀释混合能力较差、水交换缓慢。主要问题是\keyword{富营养化}。
    \item \imp{地下水}：污染过程缓慢，不易发现；自净能力极差；微生物含量少；一旦污染恢复极其困难。
    \item \imp{海洋}：污染源多而复杂；污染范围大；持续性长。
\end{enumerate}
\end{keybox}

\section{水体污染自净和污染物的转归}

\subsection{水体自净}

\begin{defbox}
\textbf{水体自净（Self-Purification）}：水体受污染后，污染物在水体的物理、化学和生物学作用下，不断稀释、扩散、分解破坏或沉入水底，水中污染物浓度逐渐降低，水质最终恢复到污染前状况的过程。
\end{defbox}

\textbf{水体自净三阶段：}
\begin{enumerate}[leftmargin=*]
    \item 第一阶段：易被氧化的有机物进行化学氧化分解——数小时完成
    \item 第二阶段：有机物在微生物作用下进行生物化学氧化分解——一般延续数日
    \item 第三阶段：含氮有机物的硝化过程——一般延续一个月左右
\end{enumerate}

\textbf{水体自净过程的特征：}污染物浓度逐渐下降；溶解氧（DO）先降后升；生物种群随之变化；水质逐渐恢复。

\textbf{环境容纳量（Environmental Capacity）}：天然水体能接纳一定量的污染物进行自净，使水质成分保持平衡的能力。

\subsection{水体自净的机制}

\begin{keybox}
\textbf{三种净化机制：}

\begin{enumerate}[leftmargin=*]
    \item \imp{物理净化}：稀释、混合、吸附、沉淀等
    \item \imp{化学净化}：氧化还原、水解、光解等
    \item \imp{生物净化}（最为重要最为活跃）：生物通过代谢作用分解水中污染物。包括：
    \begin{itemize}
        \item 需氧微生物将有机物分解为CO$_2$、水、硫酸盐、硝酸盐等
        \item 病原体的死灭
        \item 藻类对氮磷的吸收
        \item \imp{水生植物净化}：浮萍、凤眼莲、芦苇等能吸收、分解或浓缩水中汞、镉、锌等重金属及难降解有机物。例如水葫芦可吸收N、P，芦苇能分解酚类，浮萍能富集镉
    \end{itemize}
\end{enumerate}
\end{keybox}

\subsection{氧垂曲线}

\begin{keybox}
\textbf{氧垂曲线（Oxygen Sag Curve）}反映有机污染物排入河流后溶解氧的变化：

\begin{center}
\resizebox{0.95\textwidth}{!}{%
\begin{tikzpicture}[yscale=0.85, xscale=0.95]
    % Axes
    \draw[->, thick] (0,0) -- (12,0) node[below] {距离（沿河流方向）};
    \draw[->, thick] (0,0) -- (0,6.5) node[above] {溶解氧 DO (mg/L)};

    % Saturated DO line
    \draw[dashed, gray!60, thin] (0,5) -- (12,5);
    \node[gray!60, anchor=south west] at (12,5) {\footnotesize\ 饱和溶解氧};

    % DO curve
    \draw[thick, main] (0,5)
        .. controls (2,4.6) and (3,3.2) .. (4.5,2.5)
        .. controls (6,1.8) and (7.5,3.2) .. (10,4.9)
        .. controls (11,5.5) and (11.8,5.3) .. (12,5.2);
    \node[main, anchor=west] at (12,5.2) {\footnotesize\ DO曲线};

    % Critical point
    \draw[fill=keyred] (4.5,2.5) circle (2.5pt);
    \node[keyred, anchor=north west] at (4.3,2.2) {\small $C_p$（最大缺氧点）};

    % Health standard line
    \draw[red!60, dashed, thick] (0,1.5) -- (12,1.5);
    \node[red!60, anchor=south west] at (12,1.5) {\footnotesize\ 卫生标准 (4 mg/L)};

    % Zone labels
    \draw[gray!40, thick, decorate, decoration={snake, amplitude=0.5mm}] (0,5.8) -- (2.5,5.8);
    \node[gray!60] at (1.3,6.1) {\small 清洁区};

    \draw[gray!40, thick, decorate, decoration={snake, amplitude=0.5mm}] (2.5,5.8) -- (6.5,5.8);
    \node[gray!60] at (4.5,6.1) {\small 分解区（恶化区）};

    \draw[gray!40, thick, decorate, decoration={snake, amplitude=0.5mm}] (6.5,5.8) -- (12,5.8);
    \node[gray!60] at (9,6.1) {\small 恢复区};

    % Vertical dashed lines for zones
    \draw[gray!30, dashed] (2.5,0.5) -- (2.5,5.8);
    \draw[gray!30, dashed] (6.5,0.5) -- (6.5,5.8);
\end{tikzpicture}%
}
\end{center}

\begin{itemize}[leftmargin=*]
    \item \textbf{Cp点}：之前耗氧 $>$ 复氧，DO下降，水质恶化；之后复氧 $>$ 耗氧，DO逐渐恢复
    \item 若Cp点DO $>$ 4 mg/L：排放未超过水体自净能力
    \item 若Cp点DO $<$ 4 mg/L：水质严重恶化，可能出现黑臭
\end{itemize}
\end{keybox}

\subsection{水体污染物的转归}

\begin{itemize}[leftmargin=*]
    \item \imp{迁移}：污染物在水体中的空间位移
    \item \imp{生物富集作用}：生物从环境中摄入低浓度物质在体内逐渐累积
    \item \imp{生物放大作用}：污染物沿食物链逐级浓缩放大
\end{itemize}

\begin{tipbox}
\textbf{DDT在水生食物链中的迁移和转归}：水（0.00005 ppm）→ 浮游生物（0.04 ppm）→ 小鱼（0.5 ppm）→ 大鱼（2.0 ppm）→ 水鸟（25 ppm）。总体浓缩系数达50万倍。
\end{tipbox}

\section{水体污染的危害}

\subsection{生物性污染的危害}

\begin{enumerate}[leftmargin=*]
    \item \imp{介水传染病}：霍乱、伤寒、痢疾、甲型肝炎等
    \item \imp{藻类毒素污染}：
    \begin{itemize}
        \item 蓝藻水华产生\keyword{微囊藻毒素（MC）}，具有肝毒性和促癌作用。分为MC-LR（亮氨酸，急性毒性最强，抑制丝氨酸/苏氨酸蛋白磷酸酶-1和-2A活性）、MC-RR（精氨酸）、MC-YR（酪氨酸）。肝脏是MC的主要靶器官，可引起肝脏大面积肿胀出血、坏死、肝细胞结构功能破坏甚至肝衰竭
        \item MC耐热性强，煮沸不能破坏其毒性；常规水处理不能完全消除
        \item MC被确定为\imp{2B类致癌物}，是迄今已发现的最强\imp{肝癌促进剂}，且MC-LR和黄曲霉毒素B1有\imp{协同促癌作用}
        \item 慢性暴露可引起血清中谷丙转氨酶（ALT）、碱性磷酸酶（ALP）等多种酶水平显著改变
        \item 我国生活饮用水水质标准（2022）中MC-LR为扩展指标，限值1 $\mu$g/L
        \item \textbf{贝毒}：麻痹性贝毒（70\%中毒/死亡事件）、腹泻性贝毒（损伤肝脏）、记忆缺失性贝毒、神经性贝毒、西加鱼毒
    \end{itemize}
\end{enumerate}

\subsection{化学性污染的危害}

\begin{keybox}
\textbf{主要化学污染物及其危害：}

\begin{enumerate}[leftmargin=*]
    \item \imp{酚类化合物}：使水产生异臭（氯酚臭阈低至5 $\mu$g/L）；低浓度使蛋白质变性，高浓度使蛋白质沉淀；对皮肤黏膜有强烈刺激腐蚀作用；急性酚中毒表现为大量出汗、肺水肿、吞咽困难、肝及造血系统损害、黑尿
    \item \imp{多氯联苯（PCBs）}：无色或淡黄色油状有机氯化合物，性质稳定，耐酸碱、耐腐蚀。作为首批受控POPs，具持久性、蓄积性、迁移性、高毒性。主要经食物摄入暴露。\textbf{典型案例：日本米糠油中毒事件、中国台湾油症事件。}健康危害：①肝肿大和脂肪肝；②皮疹、痤疮、色素沉着，严重者肝损害、黄疸、肝性脑病；③雌激素样作用——拮抗雄激素睾酮、干扰雌激素正常代谢、干扰甲状腺功能；④通过乳汁和胎盘影响子代（胎儿油症：新生儿体重减轻、皮肤颜色异常）
    \item \imp{汞（Hg）}：转化为甲基汞后毒性大增；导致\keyword{水俣病}——慢性甲基汞中毒，以神经系统损害为主。甲基汞易于被水生生物吸收，通过生物富集和生物放大作用，使鱼贝类甲基汞富集百万倍以上
    \item \imp{邻苯二甲酸酯类（PAEs）}：塑料增塑剂、软化剂，易从塑料溶出。急性毒性很低但具有较强的雄性生殖毒性，是典型的EDCs，睾丸是重要靶器官。人类主要通过食物摄入暴露。DBP限值0.003 mg/L，DEHP限值0.008 mg/L（DEHP为2B类致癌物）
    \item \imp{全氟化合物（PFASs）}：持久性有机污染物
\end{enumerate}
\end{keybox}

\subsection{物理性污染的危害}

\begin{itemize}[leftmargin=*]
    \item \imp{热污染}：工业冷却水导致水温升高→DO降低→影响水生生态系统→促进藻类繁殖加剧富营养化。热污染还可能引起致病微生物的孳生繁殖，如变形原虫在发电厂废热水中孳生繁衍，进入人体后引起脑膜炎
    \item \imp{放射性污染}：核电站事故、核废料排放等。人体接触高浓度放射性物质的水可引起外照射，饮水或食品受放射性污染后可造成内照射
\end{itemize}

\section{水环境标准}

\subsection{地表水环境质量标准}

\textbf{制订地表水环境质量标准的原则：}
\begin{enumerate}[leftmargin=*]
    \item 防止地表水被污染而传播疾病
    \item 防止地表水被污染而危害人体健康
    \item 保证地表水感官性状良好
    \item 保证地表水自净过程正常进行
\end{enumerate}

\subsection{废水处理分级}

\begin{table}[htbp]
\centering
\caption{废水处理三级体系}
\begin{tabularx}{\textwidth}{lX}
\hline
\textbf{级别} & \textbf{处理方法与去除对象} \\
\hline
一级处理 & 物理方法——去除悬浮固体、漂浮物 \\
二级处理 & \imp{生物处理方法}——去除溶解性和胶体有机物（BOD去除） \\
三级处理 & 深度处理——去除氮磷、重金属、难降解有机物等 \\
\hline
\end{tabularx}
\end{table}

\begin{defbox}
\textbf{中水回用}：将生活污水和工业废水经处理后达到一定水质标准，回用于非饮用用途（如绿化、冲厕、工业冷却等）。
\end{defbox}

\section{水体污染的预防对策}

\begin{itemize}[leftmargin=*]
    \item \imp{污染源头控制}——污染物尚未对水体造成污染之前采用积极有效的措施，防止污染物进入水体
    \item \imp{清洁生产（Cleaner Production）}——将综合预防的环境保护策略持续应用于生产过程和产品中，以减少污染物对人类和环境的风险。内容：清洁的能源、清洁的生产过程、清洁的产品
    \item 工业废水的利用与处理
    \item 生活污水的利用与处理
    \item 医院污水的处理（含病原体、放射性物质等）
    \item 水体污染的卫生调查、监测和监督
\end{itemize}

\section{水体污染监测要点}

\subsection{江河水系监测}
\begin{enumerate}[leftmargin=*]
    \item \textbf{采样断面：}至少3个——①清洁/对照断面（污染源上游）；②污染断面（污染源下游）；③自净断面（污染断面下游至少1.5 km以上，水体基本达到自净的地方）。还需考虑表层、中层和底层的垂直布点
    \item \textbf{采样点：}水量不大/河面较窄→单点（河中心）；水量大/河面宽→多点（3--5个点，离岸50 m、150 m、江心）。采表层水应在水面下20--50 cm处
    \item \textbf{采样时间：}丰水期、枯水期和平水期各采样一次，每次连续2--3天，避开雨天。有潮汐河流应分别在高潮和低潮时采样
    \item \textbf{必测项目：}水温、浑浊度、色度、pH、总硬度、DO、BOD、总大肠菌群、挥发性酚、氰化物、砷、汞、铬。五项毒物必测：酚、氰、汞、砷、六价铬
\end{enumerate}

\subsection{底质与生物监测}
\begin{itemize}[leftmargin=*]
    \item \imp{底质监测}：底质中有害物质含量的垂直分布能反映水体污染的历史状况；某些污染物在水中含量很低不易检出，而在底质中含量可比水中高出很多倍。项目：汞、铬、砷、氰化物、酚类等五项毒物及Cu、Cd、Zn、Mn等
    \item \imp{水生生物监测}：生物种群/数量/分布测定；生物体内毒物负荷测定；污染物对水生生物综合作用检测；大肠菌群和病原微生物检测
\end{itemize}

\subsection{其他水体监测}
\begin{itemize}[leftmargin=*]
    \item \imp{湖泊水库}：增加磷（P）、氮（N）及藻类毒素的测定
    \item \imp{地下水}：增测碘、氟、砷、硫化物和硝酸盐等
\end{itemize}

\subsection{水体污染调查类型}
\begin{enumerate}[leftmargin=*]
    \item 基础调查——了解水体基本状况，范围较大
    \item 监测性调查——定期对水污染进行监测
    \item 专题调查——针对特定污染类型或污染物的专门调查
    \item 应急性突发事件调查——在严重污染事故时进行
\end{enumerate}

\section{复习题}

\begin{enumerate}[leftmargin=*, itemsep=8pt]
    \item \textbf{【名词解释】}水体自净；富营养化；水华；赤潮
    \item \textbf{【名词解释】}DO；COD；BOD$_5$；TOC；TOD
    \item \textbf{【名词解释】}硬度（暂时硬度与永久硬度）；清洁生产
    \item \textbf{【名词解释】}总大肠菌群；细菌总数
    \item \textbf{【名词解释】}氧垂曲线；生物放大作用
    \item \textbf{【简答题】}简述水体自净过程的特征及机制（含物理、化学、生物净化）。
    \item \textbf{【简答题】}简述地下水污染的特点。
    \item \textbf{【简答题】}氧垂曲线中Cp点的含义是什么？其卫生学意义何在？
    \item \textbf{【简答题】}简述水体富营养化的形成过程及危害。
    \item \textbf{【简答题】}请根据含氮化合物指标判断水体的污染状况。
    \item \textbf{【简答题】}简述制订地表水环境质量标准的原则。
    \item \textbf{【简答题】}废水处理分为哪三级？
    \item \textbf{【简答题】}简述多氯联苯（PCBs）的环境及健康危害。
    \item \textbf{【非标准答案题】}某水产养殖水体在施撒有机肥后湖面不断冒出气泡，短期内突然有大量鱼类死亡，请分析最可能的原因及应对措施。
\end{enumerate}

\subsection*{参考答案要点}

\begin{enumerate}[leftmargin=*, itemsep=5pt]
    \item \textbf{水体自净}：水体受污染后，污染物在物理、化学和生物学作用下浓度逐渐降低、水质恢复的过程。\textbf{富营养化}：湖泊水库接纳过多含P、N污水时藻类大量繁殖的现象。淡水中称\imp{水华}，海湾中称\imp{赤潮}。

    \item \textbf{DO}：溶解氧。\textbf{COD}：化学耗氧量，用强氧化剂氧化有机物消耗的氧量。\textbf{BOD$_5$}：20°C培养5日，微生物分解有机物消耗的溶解氧量。\textbf{TOC}：总有机碳。

    \item \textbf{硬度}：水中钙、镁盐类总含量（CaCO$_3$ mg/L）。暂时硬度：煮沸可去除（碳酸盐）。永久硬度：煮沸不能去除（非碳酸盐）。

    \item \textbf{总大肠菌群}：需氧及兼性厌氧、37°C使乳糖发酵产酸产气的革兰氏阴性无芽孢杆菌，粪便污染指示菌。\textbf{细菌总数}：1mL水37°C培养24h的菌落数。

    \item \textbf{氧垂曲线}：反映有机污染物排入河流后DO沿程变化的曲线。\textbf{生物放大作用}：重金属/难降解有机物通过食物链逐级浓缩放大。

    \item 特征：污染物浓度下降、DO先降后升、生物种群变化、水质恢复。机制：物理净化（稀释混合吸附沉淀）、化学净化（氧化还原水解）、生物净化（有机物降解、病原体死灭）。

    \item 污染过程缓慢不易发现；流动极缓慢自净能力差；微生物含量少；一旦污染恢复极其困难。

    \item Cp点为溶解氧最低点。之前耗氧$>$复氧，之后复氧$>$耗氧。若DO$>$4mg/L则未超自净能力；若$<$4mg/L则水质恶化。

    \item 形成：P、N污水排入→藻类大量繁殖→DO降低→水质恶化→水生生物死亡。危害：水质恶化、藻毒素产生、影响供水、破坏生态系统。

    \item NH$_3$高+NO$_3^-$低=新近污染；NH$_3$低+NO$_3^-$高=陈旧已自净；均高=持续污染。

    \item 四项原则：防止传播疾病；防止危害健康；保证感官性状良好；保证自净过程正常进行。

    \item 一级（物理去悬浮物）、二级（生物去有机物）、三级（深度处理去N、P、重金属等）。

    \item PCBs持久性、蓄积性强，具免疫、生殖、神经毒性和致癌性；可远距离迁移；米糠油事件元凶。

    \item 有机肥分解消耗大量DO→水体缺氧→鱼类窒息死亡。对策：控制施肥量、增加曝气、监测DO。
\end{enumerate}

% ----- PPT参考题 -----
\subsection*{参考题（PPT课件补充）}

\begin{enumerate}[leftmargin=*, itemsep=8pt]
    \item \textbf{【单选题】}下列水体中，卫生学特征符合"水量大、水质好"的水资源是（~~~~）
    \begin{multicols}{5}
    A. 雨雪水 \\
    B. 湖泊 \\
    C. 地下水 \\
    D. 河流 \\
    E. 海水
    \end{multicols}
    \item \textbf{【单选题】}当水中未检出亚硝酸盐氮，但氨氮和硝酸盐氮含量均明显增高时，可推测该水体受有机污染的状况是（~~~~）
    \begin{multicols}{3}
    A. 刚刚发生污染 \\
    B. 受污染不久，现仍在继续 \\
    C. 长久以来一直受污染 \\
    D. 旧污染已分解，又有新的污染 \\
    E. 已完全自净
    \end{multicols}
    \item \textbf{【单选题】}水体富营养化更容易发生在以下哪个地方？（~~~~）
    \begin{multicols}{5}
    A. 嘉陵江 \\
    B. 太平洋 \\
    C. 滇池 \\
    D. 长江 \\
    E. 亚马逊河
    \end{multicols}
\end{enumerate}

\subsubsection*{参考答案}
\begin{enumerate}[leftmargin=*, itemsep=5pt]
    \item \textbf{C（地下水）}。地下水水质清洁，矿物质含量较高，水量较稳定，符合"水量大、水质好"的特征。雨雪水水量无保证；湖泊和河流属于地表水，易受人为污染；海水为咸水，非直接可用水资源。
    \item \textbf{D（旧污染已分解，又有新的污染）}。氨氮高提示有机物正在进行氨化（新污染），硝酸盐氮高提示硝化作用已完成（旧污染已自净）。两者同时增高说明旧污染已分解且又发生了新污染。
    \item \textbf{C（滇池）}。富营养化主要发生在流动性差、水交换缓慢的封闭/半封闭水体（如湖泊、水库）。滇池属于此类湖泊，而嘉陵江、长江、亚马逊河为流动性好的河流，太平洋为开放性海域。
\end{enumerate}

\newpage

% =====================================================================
%  CHAPTER 6: 饮用水卫生
% =====================================================================
\chapter{饮用水卫生}
\setcounter{chapter}{6}
\chapterintro{
掌握介水传染病的概念及流行特点；掌握集中式供水的净化与消毒原理（混凝沉淀、过滤、氯化消毒）；掌握生活饮用水水质标准的制定原则；了解饮用水污染的应急处理原则。
}

\section{饮用水污染与疾病}

\subsection{介水传染病}

\begin{defbox}
\textbf{介水传染病（Water-borne Infection）}：由于饮用或接触受病原体污染的水而引起的一类传染病。
\end{defbox}

\begin{keybox}
\textbf{介水传染病的流行特点：}
\begin{enumerate}[leftmargin=*]
    \item 流行呈\imp{暴发型}，多数病人集中在同一潜伏期内
    \item \imp{病例分布与供水范围一致}，绝大多数患者有饮用同一水源的历史
    \item 多发生在\imp{农村地区}，饮生水者多见
    \item 一旦对污染水源采取有效措施后，疾病流行能\imp{迅速得到控制}
    \item 流行病学调查可找到水源污染途径，水质细菌学检查有异常改变
\end{enumerate}
\textbf{常见疾病：}霍乱、伤寒、副伤寒、痢疾、甲型肝炎等。
\end{keybox}

\subsection{化学性污染与健康}

\begin{enumerate}[leftmargin=*]
    \item \imp{氰化物}：剧毒，抑制细胞色素氧化酶，造成细胞内窒息。
    \begin{itemize}
        \item 毒作用机制：①CN$^-$在胃酸作用下→氰氢酸→游离氰离子+细胞色素氧化酶Fe$^{3+}$→氰化高铁细胞色素氧化酶→Fe$^{3+}$失去传递电子能力→中断呼吸链→阻断细胞内氧化代谢→细胞窒息死亡；②氰化物在体内酶作用下生成硫氰酸盐，抑制甲状腺聚碘功能，妨碍甲状腺激素合成，造成甲状腺肿大
        \item 急性中毒分4期：前驱期（刺激症状、呼气带苦杏仁味）→呼吸困难期（胸闷、心悸）→惊厥期（强制性痉挛）→麻痹期（肌肉松弛、反射消失），重者昏迷死亡
        \item 慢性中毒：神经衰弱综合征、神经细胞退行性变
        \item 饮用水中氰化物限值：0.05 mg/L
    \end{itemize}
    \item \imp{硝酸盐}：在体内可还原为亚硝酸盐，导致\keyword{高铁血红蛋白血症}（蓝婴综合征——婴幼儿特别是6个月以内婴儿血中10\%左右血红蛋白转变为高铁血红蛋白，出现发绀等缺氧症状）；与胺类生成\keyword{亚硝胺}（2A类致癌物——消化道癌症）。\textbf{婴儿为敏感人群的原因：}①婴儿胃内pH偏高（5--7）；②婴儿血红蛋白较敏感；③缺乏使高铁血红蛋白还原的酶系统；④单位体重摄水量为成年人3倍
    \item \imp{隐孢子虫}：寄生在人和动物体内，以卵囊形式从粪便排出。卵囊特点：体积小（直径4--6 $\mu$m）、感染剂量低（0.014--0.5个/L）、能抵抗常用饮水消毒剂、普通氯化消毒不能完全去除、但不耐热。\imp{贾第鞭毛虫}也是寄生于人类和动物肠道的原生动物，一般消毒方法很难将其全部杀死。两者均可通过自来水厂过滤去除
    \item \imp{饮水消毒副产物（DBPs）}：氯化消毒过程中氯与水中天然有机物（有机前体物：腐殖酸、富里酸、藻类及其代谢物、蛋白质等）反应生成。包括三卤甲烷（THMs）、卤乙酸（HAAs）等，具有"三致"效应。影响因素：有机前体物含量、加氯量、溴离子浓度、pH（THMs随pH↑而↑，HAAs随pH↑而↓）、温度、接触时间。"Mutagen X"（MX，3-氯-4-二氯甲基-5-羟基-2(5H)呋喃酮）在氯化消毒饮水中浓度很低（数ng/L至数十ng/L），却是\imp{迄今为止最强的诱变物之一}。NDMA（N-二甲基亚硝胺）为2A类致癌物。流行病学研究在膀胱癌和直肠癌显示与DBPs有一致性
\end{enumerate}

\subsection{二次供水污染}

\begin{defbox}
\textbf{二次供水（Secondary Water Supply）}：用水单位将来自市政供水或自备井的生活饮用水贮存于储水池或水箱中，再通过机械加压或凭借高层建筑的自然压差，二次输送至用户的供水系统。我国高层建筑5层以上用水都需要二次供水系统。
\end{defbox}

\begin{tipbox}
\textbf{高层建筑二次供水污染的原因：}
\begin{itemize}[leftmargin=*]
    \item 储水设备问题：储水箱材料不合格（材质老化、侵蚀、涂料脱落）；水箱容积过大，水滞留时间过长，余氯耗尽，微生物繁殖；储水箱设计不合理（出水口高出水箱底平面形成死水）；未及时维修保养、未定期清洗消毒
    \item 管网问题：管道错误连接（上水管设在污水管下面、溢水管与污水管相连无防倒灌措施）；管网破裂渗漏、降压和停水维修时外界污物渗入；管网陈旧老化，水管材质和涂料脱落析出
\end{itemize}
\textbf{健康危害：}以生物性污染为主，通常引起介水肠道传染病；若储水箱材料中有害物质（铅、镉、砷、汞等）含量或溶出过高，可导致慢性危害。
\end{tipbox}

\section{生活饮用水水质标准}

\begin{keybox}
\textbf{制定原则：}
\begin{enumerate}[leftmargin=*]
    \item \imp{流行病学上安全}——不含病原微生物
    \item \imp{化学组成对人体无害}——不引起急慢性中毒及远期危害
    \item \imp{感官性状良好}——色、臭、味、浑浊度满足要求
    \item \imp{经济技术上可行}——标准限值可在现实条件下实现
\end{enumerate}
\end{keybox}

常规检测项目：感官性状和一般化学指标、毒理学指标、细菌学指标、放射性指标。

\begin{tipbox}
\textbf{生活饮用水水质常规指标分类（GB 5749-2022）：}
\begin{itemize}[leftmargin=*]
    \item \imp{微生物学指标}：总大肠菌群（不得检出/100mL）、耐热大肠菌群（不得检出/100mL）、大肠埃希菌（不得检出/100mL）、菌落总数（$\leq$100 CFU/mL）。\textbf{耐热大肠菌群/粪大肠菌群}是判断饮用水是否受粪便污染的重要指标
    \item \imp{消毒剂指标}：游离氯出厂水$\geq$0.3 mg/L（接触30min），末梢水$\geq$0.05 mg/L；末梢水游离氯余量可作为输配水过程中有无再次污染的信号
    \item \imp{毒理学指标}：砷（I类致癌物）、镉、铬（六价）、铅、汞、硒、氰化物、氟化物（$\leq$1 mg/L）、硝酸盐、三氯甲烷、四氯化碳、溴酸盐、甲醛、亚氯酸盐、氯酸盐
    \item \imp{放射性指标}：总$\alpha$放射性、总$\beta$放射性
    \item \imp{感官性状和一般化学指标}：色度、浑浊度（限值1 NTU）、臭和味、肉眼可见物、pH（6.5--8.5）、铝、铁（$\leq$0.3 mg/L）、锰（$\leq$0.1 mg/L）、铜、锌、氯化物、硫酸盐、溶解性总固体、总硬度（$\leq$450 mg/L）、耗氧量、挥发酚类、阴离子合成洗涤剂
\end{itemize}

\textbf{非常规指标：}微生物学2项（贾第鞭毛虫、隐孢子虫）；毒理学39项（农药、除草剂、苯化合物、微囊藻毒素-LR等）；感官性状3项。
\end{tipbox}

\begin{tipbox}
\textbf{部分水质标准制定依据：}
\begin{itemize}[leftmargin=*]
    \item \imp{浑浊度}：限值1度。10度时人即可觉察浑浊；降低浑浊度对去除有害物质、提高消毒效果有积极作用
    \item \imp{pH}：限值6.5--8.5。过低腐蚀管道，过高析出溶解性盐类并降低氯消毒效果
    \item \imp{总硬度}：限值450 mg/L（以CaCO$_3$计）。硬度过高可引起胃肠功能暂时性紊乱
    \item \imp{铁}：限值0.3 mg/L。1 mg/L时有明显金属味，0.5 mg/L时色度可$>$30度
    \item \imp{锰}：限值0.1 mg/L。微量即可使水呈黄褐色；超过0.15 mg/L时使衣物和瓷器出现色斑。锰虽有神经毒性，但因感官影响阈值低于毒理学阈值，按最敏感指标原则归为感官性状指标
    \item \imp{氟化物}：限值1 mg/L。综合考虑1 mg/L时对牙齿的轻度影响和防龋作用，及高氟地区除氟在经济技术上的可行性
\end{itemize}
\end{tipbox}

\section{集中式给水}

\subsection{水源选择原则}

\begin{enumerate}[leftmargin=*]
    \item \imp{水量充足} \item \imp{水质良好} \item \imp{便于防护} \item \imp{技术经济合理}
\end{enumerate}

\subsection{水质净化与消毒}

\begin{keybox}
\textbf{（一）混凝沉淀（Coagulation Precipitation）}

\textbf{原理：}
\begin{enumerate}[leftmargin=*]
    \item \imp{压缩双电层作用}——降低$\zeta$电位
    \item \imp{电中和作用}——中和胶体负电荷
    \item \imp{吸附架桥作用}——高分子混凝剂链状结构吸附多个胶体颗粒
\end{enumerate}

\textbf{混凝剂比较：}

\begin{table}[H]
\centering
\caption{常用混凝剂比较}
\begin{tabularx}{\textwidth}{lXX}
\hline
\textbf{特性} & \textbf{铝盐（明矾、硫酸铝）} & \textbf{铁盐（三氯化铁、硫酸亚铁）} \\
\hline
适宜pH & 窄（6.5--7.5） & 宽（5--9） \\
絮凝体 & 慢且松散（低温时更明显） & 快、大而紧密 \\
腐蚀性 & 小 & 大，易潮湿 \\
对水质影响 & 无不良影响 & 处理后含铁量高 \\
低温低浊水效果 & 较差 & 较好 \\
\hline
\end{tabularx}
\end{table}

\imp{聚合氯化铝（PAC）}：适用pH范围宽、腐蚀性小、成本低、絮体形成快、沉淀效果好、用量少。

\imp{聚丙烯酰胺（PAM）}：高分子混凝剂，常作为助凝剂使用。

\textbf{助凝剂}：①调节混凝条件（原水碱度不够加石灰；用氯将亚铁氧化为高铁）；②改善絮凝体结构（铝盐产生的絮凝体细小松散时，用聚丙烯酰胺或活性硅酸等助凝）。

\textbf{影响因素：}水温；水的pH（对金属盐类混凝剂影响较大，对高分子混凝剂影响较小）；水中微粒的性质和含量；水中有机物和溶解盐含量；混凝剂种类和用量；投加方法、搅拌强度和反应时间等。

\textbf{（二）过滤（Filtration）}

\textbf{原理：}浑水通过滤料层，悬浮颗粒和微生物被截留（筛除作用+接触凝聚作用）。工作周期：成熟期→过滤期→清洗期。

\textbf{滤料要求：}本身无毒、化学性质稳定；不能被微生物利用和分解；有良好机械强度；颗粒粒度均匀，有适当孔隙率。常用滤料：石英砂、无烟煤、木炭、活性炭、磁铁砂。

\textbf{影响因素：}滤层厚度和粒径、滤速和滤池类型、进水水质（浑浊度、色度、有机物、藻类）。

\textbf{（三）消毒（Disinfection）}

\textbf{理想消毒剂的选择标准：}①能迅速灭活病原微生物（细菌+病毒+真菌+寄生虫）；②在使用的浓度水平无毒；③便于控制和检测；④有剩余消毒剂；⑤对水感官性状无不良影响；⑥副产物对健康影响小且可预防或消除；⑦经济技术上可行。

\textbf{氯化消毒（Chlorination Disinfection）}

\textbf{原理：}含氯制剂加入水中生成\imp{次氯酸（HOCl）}，HOCl为中性小分子，强渗透性，强氧化剂，损害细胞膜使通透性增加，蛋白质、RNA、DNA等释放，影响多种酶系统（细菌：干扰糖代谢；病毒：对核酸致死性损害）。

\begin{defbox}
\textbf{有效氯}：含氯化合物分子团中氯的价数 $>$ -1，是含氯化合物中有杀菌能力的有效成分。

\textbf{需氯量}：因杀菌、氧化水中有机物和还原性无机物需要的总氯量。

\textbf{余氯（Residual Chlorine）}：加氯量 $>$ 需氯量，使在氧化和杀菌后还能剩余的有效氯量。是评价氯化消毒效果是否达标的简便指标。
\begin{itemize}
    \item \imp{游离性余氯}：HOCl + OCl$^-$，消毒效果好。接触30 min，含量0.3--0.5 mg/L
    \item \imp{化合性余氯}：NH$_2$Cl + NHCl$_2$，消毒效果略差。接触1--2 h，含量1--2 mg/L
\end{itemize}
\end{defbox}

\textbf{HOCl vs OCl$^-$}：HOCl的杀菌效率比OCl$^-$高80倍。pH $<$ 6时HOCl $>$ 95\%；pH = 7.5时HOCl = OCl$^-$；pH $>$ 9时OCl$^-$ $\approx$ 100\%。

\textbf{影响因素：}
\begin{enumerate}[leftmargin=*]
    \item \imp{加氯量}——越多效果越好，但易产生DBPs
    \item \imp{接触时间}——越长效果越好。水温每提高10°C，杀菌率提高2--3倍
    \item \imp{pH值}——pH越低HOCl比例越高，效果越好
    \item \imp{水温}——越高杀菌效果越好 → 冬季需延长接触时间
    \item \imp{浑浊度}——高时细菌附着在悬浮颗粒上，氯不易直接作用于细菌
    \item \imp{微生物种类和数量}——抵抗力：大肠杆菌 $<$ 病毒 $<$ 原虫包囊
\end{enumerate}

\textbf{氯化消毒方法：}
\begin{itemize}[leftmargin=*]
    \item \imp{普通氯化消毒法}：水浊度低、有机物污染轻、基本无氨（$<$0.3 mg/L）无酚时使用。接触时间短，效果可靠，但产生氯化副产物
    \item \imp{氯胺消毒法}：水中加氯后再加氨，生成一氯胺和二氯胺。副产物少，稳定性高，适宜长距离输水；但消毒作用不如次氯酸强，接触时间长
    \item \imp{折点氯消毒法}：加氯量超过折点，形成适量游离氯。消毒效果可靠，能降低臭味、色度及有机物含量；但耗氯多，副产物多
    \item \imp{过量氯消毒法}：水源污染严重或需短时间消毒时，加过量氯使余氯达1--5 mg/L，消毒后用亚硫酸钠/活性炭脱氯
\end{itemize}
\end{keybox}

\subsection{消毒方法比较}

\begin{table}[htbp]
\centering
\caption{常用饮水消毒方法比较}
\begin{tabularx}{\textwidth}{lXXXX}
\hline
\textbf{特性} & \textbf{氯化消毒} & \textbf{二氧化氯} & \textbf{臭氧消毒} & \textbf{紫外线消毒} \\
\hline
消毒效果 & 良好 & 优异 & 优异 & 良好 \\
持续消毒能力 & \imp{有} & 有 & 无 & 无 \\
消毒副产物 & THMs等 & 亚氯酸盐 & 溴酸盐等 & 无 \\
成本 & 低 & 中等 & 高 & 中等 \\
\hline
\end{tabularx}
\end{table}

\subsection{氯化消毒副产物的预防}

\begin{itemize}[leftmargin=*]
    \item 减少水源中有机物前体物质
    \item 优化消毒工艺（控制加氯量、改变加氯点）
    \item 采用替代消毒剂（氯胺、二氧化氯、臭氧）
    \item 深度处理去除已生成的副产物
\end{itemize}

\subsection{饮用水的深度净化}

在市政供水常规处理基础上，再次对水质净化处理以获得优质饮用水。
\begin{itemize}[leftmargin=*]
    \item \imp{物理吸附分离}：活性炭吸附法；膜分离法（微滤MF、超滤UF、纳滤NF、反渗透RO）
    \item \imp{化学氧化技术}：预氧化技术、TiO$_2$光催化技术
    \item \imp{生物预处理技术}
\end{itemize}

\section{分散式给水}

\subsection{包装饮用水}
\begin{itemize}[leftmargin=*]
    \item \imp{纯水}：以市政自来水为原水，经反渗透、电渗析、蒸馏等工艺使水中溶解矿物质及其他有害物质全部去除，电导率$<$10 $\mu$s/cm，浊度$<$1度。\textbf{不建议作生活饮用水长期饮用}
    \item \imp{净水}：以市政自来水为原水，通过吸附、超滤、纳滤去除有害物质而保留原水中溶解性矿物质。\textbf{可作生活饮用水长期饮用}
    \item \imp{天然矿泉水}：储存于地下深处自然涌出或人工采集的未受污染且含偏硅酸、锶、锌等一种或多种微量元素达到限量值的泉水
\end{itemize}

\subsection{直饮水（Direct Drinking Water）}
属分质供水范畴，对自来水深度处理后，通过食品卫生级输水管道送入用户，可直接饮用。

\subsection{涉水产品}
在饮用水生产和供水过程中与饮用水接触的连接止水材料、塑料及有机合成管材、管件、防护涂料、水处理剂、水质处理器及其他新材料和化学物质。

\subsection{传统分散式给水}

\begin{itemize}[leftmargin=*]
    \item \textbf{地下水}（井水）：水质较好，需注意卫生防护
    \item \textbf{地表水}：需经处理后方可饮用
    \item \textbf{雨雪水}：水量无保证，易受大气污染影响
\end{itemize}

\section{饮用水污染的应急处理}

\textbf{水污染事故应急处理原则：}
\begin{enumerate}[leftmargin=*]
    \item 立即停止使用受污染水源
    \item 迅速查明污染源和污染物
    \item 采取紧急净化和消毒措施
    \item 启用备用水源或应急供水
    \item 开展人群健康监测和医学救援
    \item 及时向公众通报信息
\end{enumerate}

\section{复习题}

\begin{enumerate}[leftmargin=*, itemsep=8pt]
    \item \textbf{【名词解释】}介水传染病；饮水消毒副产物（DBPs）
    \item \textbf{【名词解释】}混凝沉淀；氯化消毒
    \item \textbf{【简答题】}简述介水传染病的流行特点。
    \item \textbf{【简答题】}简述生活饮用水水质标准的制定原则。
    \item \textbf{【简答题】}简述集中式供水水源选择的原则。
    \item \textbf{【简答题】}简述混凝沉淀的原理及影响因素。
    \item \textbf{【简答题】}简述影响氯化消毒效果的因素。
    \item \textbf{【简答题】}比较四种常用消毒方法的优缺点。
    \item \textbf{【简答题】}简述如何预防和控制氯化消毒副产物。
    \item \textbf{【简答题】}简述高层建筑二次供水污染的原因。
    \item \textbf{【简答题】}饮用水应急事件的调查内容和处理原则是什么？
    \item \textbf{【非标准答案题】}理想的饮用水应该是什么样的？为什么长期饮用纯净水对健康可能不利？
\end{enumerate}

\subsection*{参考答案要点}

\begin{enumerate}[leftmargin=*, itemsep=5pt]
    \item \textbf{介水传染病}：通过饮用或接触受病原体污染的水而引起的传染病。\textbf{DBPs}：消毒剂与原水中有机物反应生成的具有"三致"效应的副产物。

    \item \textbf{混凝沉淀}：加入混凝剂使水中悬浮物和胶体聚合沉淀的净水方法。\textbf{氯化消毒}：使用含氯制剂生成HOCl杀菌的消毒方法。

    \item 五大特点：暴发型、病例分布与供水范围一致、多见于农村、控制水源后迅速控制、水质细菌学异常。

    \item 四项原则：流行病学安全、化学物质无害、感官性状良好、经济技术可行。

    \item 四项原则：水量充足、水质良好、便于防护、技术经济合理。

    \item 原理：压缩双电层+电中和+吸附架桥。影响因素：pH、水温、混凝剂种类用量、搅拌。

    \item 六大因素：加氯量、接触时间、pH值、水温、浑浊度、微生物种类和数量。

    \item 氯化持续效果好成本低但有副产物；臭氧效果好但无持续能力成本高；紫外线无副产物但无持续能力；二氧化氯效果好有持续能力副产物较少。

    \item 减少前体物质、优化消毒工艺、采用替代消毒剂、深度处理去除已生成DBPs。

    \item 储水池设计不合理、储水时间过长、管理不善、管网老化渗漏。

    \item 调查：污染源、污染物、范围、暴露人群、健康影响。处理：停用、查明、紧急处理、备用水源、健康监测、信息公开。

    \item 理想饮用水应安全卫生、含适量矿物质、感官良好。纯净水缺乏Ca、Mg等矿物质，pH偏酸，长期饮用影响矿物质平衡。
\end{enumerate}

% ----- PPT参考题 -----
\subsection*{参考题（PPT课件补充）}

\begin{enumerate}[leftmargin=*, itemsep=8pt]
    \item \textbf{【单选题】}下列哪种疾病属于介水传染病？（~~~~）
    \begin{multicols}{5}
    A. 台湾油症 \\
    B. 霍乱 \\
    C. 水俣病 \\
    D. SARS \\
    E. 氟骨症
    \end{multicols}
    \item \textbf{【单选题】}针对介水传染病的暴发流行，最有效的防控措施是（~~~~）
    \begin{multicols}{5}
    A. 勤洗手 \\
    B. 戴防护口罩 \\
    C. 隔离感染病人 \\
    D. 饮水净化消毒 \\
    E. 接种疫苗
    \end{multicols}
    \item \textbf{【单选题】}有村民反映当地饮用水烧开后结垢多，CDC取水样检测发现水中溶解性总固体为1872 mg/L，总硬度550 mg/L，大大超过生活饮用水卫生标准。这样的水最可能来自（~~~~）
    \begin{multicols}{5}
    A. 江水 \\
    B. 河水 \\
    C. 湖泊水 \\
    D. 井水 \\
    E. 水库水
    \end{multicols}
    \item \textbf{【单选题】}锰具有神经毒性，微量锰可使水呈黄褐色。在生活饮用水卫生标准中，锰被归类为感官性状和一般化学指标，而不是毒理学指标，这是遵循了环境卫生标准制订中的（~~~~）
    \begin{multicols}{3}
    A. 保障居民不发生急性及慢性危害原则 \\
    B. 最敏感指标的原则 \\
    C. 对人体健康无间接影响原则 \\
    D. 最经济指标的原则 \\
    E. 最切实可行的原则
    \end{multicols}
    \item \textbf{【单选题】}卫生监督执法部门在对某自来水厂进行常规监督时发现，该水厂出厂水游离氯余量为0.5 mg/L，而末梢水中几乎不能检测到游离余氯，提示（~~~~）
    \begin{multicols}{2}
    A. 水厂液氯添加量不足 \\
    B. 自来水厂操作不规范 \\
    C. 管网内出现二次污染 \\
    D. 游离余氯转变为化合性余氯
    \end{multicols}
\end{enumerate}

\subsubsection*{参考答案}
\begin{enumerate}[leftmargin=*, itemsep=5pt]
    \item \textbf{B（霍乱）}。介水传染病指通过饮用或接触受病原体污染的水而引起的传染病，霍乱是典型的介水传染病。台湾油症由多氯联苯引起（食物污染），水俣病为慢性甲基汞中毒（化学性污染），SARS为呼吸道传染病，氟骨症为地球化学性疾病。
    \item \textbf{D（饮水净化消毒）}。介水传染病的传播途径是受污染的水源，最有效的控制措施是切断传播途径——对饮用水进行净化消毒。勤洗手属个人卫生，口罩防呼吸道传播，隔离病人针对传染源，疫苗接种针对易感人群，均不如饮水消毒直接有效。
    \item \textbf{D（井水）}。地下水（井水）溶解土壤矿物质较多，硬度较高，溶解性总固体含量高于地表水。江、河、湖、水库水均属地表水，硬度通常较低。
    \item \textbf{B（最敏感指标的原则）}。锰在低于产生神经毒性的浓度时即可使水着色（感官不良），其感官影响的阈值低于毒理学阈值，按最敏感指标的原则归类为感官性状和一般化学指标。
    \item \textbf{C（管网内出现二次污染）}。出厂水有余氯（0.5 mg/L）说明水厂加氯正常；末梢水无游离余氯说明余氯在管网输送途中被消耗，最常见的原因是管网老化渗漏导致二次污染，细菌等污染物消耗了余氯。
\end{enumerate}

\newpage

% =====================================================================
%  CHAPTER 7: 土壤卫生 【重点章节】
% =====================================================================
\chapter{土壤卫生  【重点章节】}
\setcounter{chapter}{7}
\chapterintro{
本章为\keyword{重点考察章节}（3学时）。掌握土壤环境背景值的概念及意义；掌握土壤污染的基本特点；掌握土壤重金属污染（镉、铊、铬等）和农药污染的健康危害；掌握持久性有机污染物（POPs）的特性及危害；了解土壤卫生防护措施。
}

\section{土壤环境特征}

\subsection{土壤的组成}
\begin{itemize}[leftmargin=*]
    \item \imp{土壤固相}：矿物质（占土壤固体总质量90\%以上）+ 有机物质（腐殖质、生物残体、土壤生物）
    \item \imp{土壤液相}：水分及其水溶物质
    \item \imp{土壤气相}：土壤空隙中多种气体混合物；上层与大气相似，深层O$_2$减少、CO$_2$增加
\end{itemize}

\subsection{土壤物理学特征}

\begin{defbox}
\textbf{土壤质地}：按土壤粒级及其组合比例而定的土壤名称，反映土壤固相颗粒分布情况。分为\imp{砂土、壤土、黏土}。壤土通气透水、保水保肥性能均较好，是良好的土壤质地。

\textbf{土壤孔隙度（Soil Porosity）}：自然状态下，单位容积土壤中孔隙容积所占百分率。对土壤性质的影响：①容水性（保存水分的能力）；②渗水性（水分渗透的能力）；③透气性（空气穿透的能力）；④毛细管现象（水分沿孔隙上升的程度）。

\textbf{规律：}土壤颗粒越小，容水性越大，透气性越差。
\end{defbox}

\begin{keybox}
\textbf{土壤环境背景值（Background Level）}：该地区未受或少受人为污染影响的天然土壤中各种元素的含量。

\textbf{实践意义：}
\begin{enumerate}[leftmargin=*]
    \item 评价化学污染物对土壤污染程度的\imp{参照值}
    \item 制定土壤中有害化学物质卫生标准的\imp{重要依据}
    \item 评价土壤化学环境对居民健康影响的\imp{重要依据}
    \item 土地资源开发利用和\imp{地方病防治}的科学依据
\end{enumerate}
\end{keybox}

\begin{defbox}
\textbf{土壤环境容量/土壤负载容量}：一定土壤环境单元在一定时限内遵循环境质量标准，即维持土壤生态系统的正常结构与功能，保证农产品的生物学产量与质量，在不使环境系统污染的前提下，土壤环境所能容纳污染物的最大负荷量。是充分利用土壤环境的纳污能力、实现污染物总量控制、合理制定环境质量标准和卫生标准的重要依据。

\textbf{腐殖质（Humus）}：动植物残体经分解转化形成的有机物质，占土壤有机质总量的\imp{85\%--90\%}。
\end{defbox}

\begin{tipbox}
\textbf{土壤生物性污染监测指标：}
\begin{itemize}[leftmargin=*]
    \item \imp{大肠菌值}：发现大肠菌的最少土壤克数，代表人畜粪便污染的主要指标，也是代表肠道传染病危险性的主要指标
    \item \imp{产气荚膜杆菌值}：代表粪便污染的指标，与大肠菌在土壤中数量的消长关系可判定污染时间长短——产气荚膜杆菌多而大肠菌相对少→\textbf{陈旧性污染}；大肠菌值更多→\textbf{新鲜污染}（危害性较大）
    \item \imp{蛔虫卵数}：直接说明在流行病学上是否对人体健康有威胁。根据蛔虫卵在土壤中的不同发育阶段及活卵所占百分比，可判断土壤的自净程度
\end{itemize}
\end{tipbox}

\section{土壤的污染、自净及转归}

\subsection{土壤污染的基本特点}

\begin{keybox}
\textbf{四大基本特点：}
\begin{enumerate}[leftmargin=*]
    \item \imp{隐蔽性}——不易被直接察觉，往往通过食物链传递才被发现
    \item \imp{累积性与地域性}——污染物长期累积，区域差异大
    \item \imp{不可逆转性}——重金属等一旦进入极难清除
    \item \imp{治理的长期性}——修复周期长、成本高
\end{enumerate}
\end{keybox}

\subsection{污染方式}

\begin{enumerate}[leftmargin=*]
    \item \imp{气型污染}——大气污染物沉降
    \item \imp{水型污染}——工业废水和生活污水通过污水灌溉污染
    \item \imp{固体废弃物型污染}——垃圾、工业废渣等
\end{enumerate}

\subsection{土壤自净}

\begin{defbox}
\textbf{土壤自净作用（Soil Self-Purification）}：受污染土壤通过物理、化学和生物学作用，使病原体死灭，有害物质转化为无害，土壤逐渐恢复到污染前状态的过程。
\end{defbox}

\textbf{含氮有机物的无机化过程：}蛋白质等 → 氨基酸 → NH$_3$（\imp{氨化作用}）→ NO$_2^-$（亚硝化作用）→ NO$_3^-$（\imp{硝化作用}）

\subsection{重金属在土壤中的转化}

\begin{itemize}[leftmargin=*]
    \item 土壤胶体、腐殖质的吸附和螯合作用
    \item \imp{土壤pH}：pH越低，重金属溶解度越高，易被植物吸收
    \item 土壤氧化还原状态的影响
\end{itemize}

\section{土壤污染与健康}

\subsection{重金属污染}

\begin{keybox}
\textbf{镉（Cd）污染——痛痛病（Itai-itai Disease）：}
\begin{itemize}[leftmargin=*]
    \item 来源：矿山废水、冶炼、磷肥。水体镉污染主要来自工业废水；土壤镉污染主要是含镉废水灌溉农田所致
    \item IARC已将镉列为\imp{I类致癌物}。镉在人体中生物半减期长达\imp{10--25年}，可在体内不断累积。肾脏可蓄积镉吸收量的1/3，是镉中毒的主要靶器官
    \item 机制：镉与含羟基、氨基、巯基的蛋白质分子结合→抑制酶系统→影响肝肾功能；镉损伤肾小管→维生素D$_3$活性受影响→骨骼生长代谢受阻→钙磷代谢障碍→骨质疏松软化；镉干扰铁代谢→肠道铁吸收减低、破坏红细胞→贫血
    \item 表现：全身剧烈刺痛（腰、背、膝关节→遍及全身），骨骼严重畸形，骨脆易折，多发性病理骨折，肾功能损害。尿镉反映镉性肾损伤和体内镉负荷（以5 $\mu$mol/mol肌酐为慢性镉中毒诊断下限值），血镉反映近期接触量
    \item 流行病学：多为40岁以上妇女，妊娠、哺乳、营养不良等是发病诱因，常分布在镉矿丰富、雨量大、坡度陡、水土流失严重地区
\end{itemize}

\textbf{铊（Tl）污染：}神经系统和消化系统损害，特征性脱发

\textbf{铬（Cr）污染：}Cr(VI)毒性远大于Cr(III)；皮肤溃疡（铬疮）、呼吸道刺激、致癌性
\end{keybox}

\subsection{农药污染与POPs}

\begin{keybox}
\textbf{持久性有机污染物（POPs）的四大特性：}
\begin{enumerate}[leftmargin=*]
    \item \imp{持久性}——在环境中难以降解
    \item \imp{蓄积性}——易在生物体内脂肪组织中蓄积
    \item \imp{迁移性}——可通过大气、水流远距离迁移（"全球蒸馏效应"）
    \item \imp{高毒性}——具有致癌、致畸、致突变"三致"效应
\end{enumerate}

\textbf{健康危害：}急慢性毒性；"三致"作用；干扰生殖内分泌功能；免疫毒性、神经发育毒性。
\end{keybox}

\subsection{生物性污染的危害}

\begin{enumerate}[leftmargin=*]
    \item \imp{肠道传染病和寄生虫病}——经口传播
    \item \imp{钩端螺旋体病和炭疽病}——接触污染土壤或水
    \item \imp{破伤风和肉毒中毒}——伤口接触含芽孢的土壤
\end{enumerate}

\section{土壤污染的预防对策}

\begin{itemize}[leftmargin=*]
    \item \imp{土壤污染防治法律法规}
    \item \imp{土壤环境质量标准}
    \item \imp{土壤卫生防护}：粪便无害化处理、垃圾无害化处理（卫生填埋、焚烧、堆肥）、危险工业废渣处理、污水灌溉卫生防护
    \item \imp{污染土壤修复}：物理修复、化学修复、生物修复
\end{itemize}

\section{复习题}

\begin{enumerate}[leftmargin=*, itemsep=8pt]
    \item \textbf{【名词解释】}土壤环境背景值；环境容量；腐殖质
    \item \textbf{【名词解释】}土壤污染；土壤自净；POPs
    \item \textbf{【简答题】}简述土壤环境背景值的实践意义。
    \item \textbf{【简答题】}简述土壤污染的基本特点。
    \item \textbf{【简答题】}简述土壤自净的机制（以含氮有机物无机化为例）。
    \item \textbf{【简答题】}试述镉污染与痛痛病的关系（来源、机制、表现）。
    \item \textbf{【简答题】}POPs有哪些特性？对健康有何危害？
    \item \textbf{【简答题】}简述土壤生物性污染引起的主要疾病及危害途径。
    \item \textbf{【简答题】}简述土壤的卫生防护原则。
    \item \textbf{【非标准答案题】}结合痛痛病和水俣病的案例，谈谈食物链传递在环境污染健康危害中的关键作用。
\end{enumerate}

\subsection*{参考答案要点}

\begin{enumerate}[leftmargin=*, itemsep=5pt]
    \item \textbf{背景值}：未受污染天然土壤中各元素含量。\textbf{环境容量}：不超过卫生标准前提下污染物能容纳的最大负荷量。\textbf{腐殖质}：动植物残体经分解转化形成的有机物质。

    \item \textbf{土壤污染}：人类活动排放有害物质进入土壤危害人畜健康。\textbf{土壤自净}：受污染土壤通过物化生作用恢复的过程。\textbf{POPs}：具持久性、蓄积性、迁移性和高毒性的有机污染物。

    \item 四项意义：评价污染程度的参照值；制定卫生标准的依据；评价对健康影响的依据；土地开发和地方病防治的科学依据。

    \item 四大特点：隐蔽性、累积性与地域性、不可逆转性、治理长期性。

    \item 物理+化学+生物净化。含氮有机物无机化：蛋白质→氨基酸→NH$_3$（氨化）→NO$_2^-$（亚硝化）→NO$_3^-$（硝化）。

    \item 来源：矿山废水、冶炼、磷肥。机制：Cd蓄积→肾小管损伤→钙磷障碍→骨软化。表现：全身剧烈疼痛、骨折、肾功能衰竭。

    \item 四特性：持久性、蓄积性、迁移性、高毒性。危害：急慢性毒性、三致效应、干扰内分泌、免疫和神经毒性。

    \item 肠道传染病和寄生虫病（经口）、钩端螺旋体病和炭疽病（接触）、破伤风和肉毒中毒（伤口感染）。

    \item 法律标准制定、土壤卫生防护（粪便垃圾无害化等）、发展生态农业、污染土壤修复。

    \item 围绕食物链的生物放大效应、从环境到人体的传递路径、对风险预防和治理的启示展开。
\end{enumerate}

% ----- PPT参考题 -----
\subsection*{参考题（PPT课件补充）}

\subsubsection*{一、思考题}

\begin{enumerate}[leftmargin=*, itemsep=8pt]
    \item \textbf{【思考题】}某一地区农用土壤样品检测，11项指标（8种重金属+2种农药+苯并[a]芘）均低于风险筛选值，是否表明该土壤一定无污染？请说明理由。
\end{enumerate}

\subsubsection*{二、判断题（判断下列说法是否正确，错误的请说明理由）}

\begin{enumerate}[leftmargin=*, itemsep=8pt]
    \item 土壤污染引起破伤风的主要危害途径是土壤-人。（~~~~）
    \item 土壤中元素的本底值是指任何一处土壤中各种元素的含量。（~~~~）
    \item 土壤的pH较高，呈碱性时有利于农作物吸收重金属元素。（~~~~）
    \item 土壤污染的基本特点包括隐蔽性、累积性、迁移性、长期性。（~~~~）
    \item 持久性有机污染物特性包括持久性、迁移性、生物蓄积性、高毒性。（~~~~）
    \item 镉污染所致痛痛病属于土壤水型污染。（~~~~）
    \item 土壤污染的来源包括地质环境中区域性差异导致土壤中某些元素过高。（~~~~）
    \item 土壤颗粒越小，土壤容水性越大，透气性越差。（~~~~）
    \item 土壤气型污染的特点之一是污染物集中在土壤深层。（~~~~）
    \item 镉、汞、砷、铅、铬含量高于风险管制值的土地，原则上禁止种植食用农产品。（~~~~）
\end{enumerate}

\subsubsection*{参考答案}

\textbf{思考题：}
\begin{enumerate}[leftmargin=*, itemsep=5pt]
    \item 不一定。理由：①我国现行农用地土壤标准（GB 15618-2018）仅涵盖11项指标，但土壤污染物种类远多于此；②例如铊（Tl）已被加拿大、德国、美国等国家列为优先控制污染物，但我国尚未将其纳入土壤环境质量标准；③可能存在"假阴性"——未纳入标准的污染物即使超标也无法检出。因此，应关注标准外的潜在污染物。
\end{enumerate}

\textbf{判断题：}
\begin{enumerate}[leftmargin=*, itemsep=5pt]
    \item \textbf{×（错误）}。破伤风的主要危害途径是\textbf{土壤→伤口→人}，即含破伤风芽孢的土壤通过破损皮肤进入人体，而非直接"土壤-人"。
    \item \textbf{×（错误）}。本底值（背景值）是指\textbf{未受人为污染的天然土壤}中各种元素的含量，而非"任何一处土壤"。受污染地区的土壤元素含量不属于本底值。
    \item \textbf{×（错误）}。土壤pH越高（碱性越强），重金属溶解度越低，形成不溶性氢氧化物，反而\textbf{不利于}农作物吸收。酸性土壤中重金属活性更高，更易被作物吸收。
    \item \textbf{×（错误）}。土壤污染四大基本特点为：\textbf{隐蔽性、累积性与地域性、不可逆转性、治理长期性}。"迁移性"不属于基本特点（迁移性是POPs的特性）。
    \item \textbf{√（正确）}。POPs的四大特性为持久性、迁移性（可远距离传输）、生物蓄积性、高毒性。
    \item \textbf{√（正确）}。痛痛病由镉污染通过废水→土壤→稻米→人体的水型污染途径引起。
    \item \textbf{×（错误）}。地质环境中区域性差异导致土壤中某些元素过高属于\textbf{原生环境}问题（如生物地球化学性疾病），不属于人为活动造成的"土壤污染"。
    \item \textbf{√（正确）}。土壤颗粒越小，孔隙越小，容水性（持水性）越大，但透气性越差。
    \item \textbf{×（错误）}。气型污染主要由大气污染物沉降引起，污染物主要集中在土壤\textbf{表层}，而非深层。
    \item \textbf{√（正确）}。根据《农用地土壤污染风险管控标准》，重金属含量高于风险管制值的土地，原则上应采取禁止种植食用农产品等严格管控措施。
\end{enumerate}

\newpage

% =====================================================================
%  CHAPTER 8: 环境相关疾病
% =====================================================================
\chapter{环境相关疾病}
\setcounter{chapter}{8}
\chapterintro{
掌握生物地球化学性疾病的概念及流行特征；掌握碘缺乏病、地方性氟中毒、地方性砷中毒的病因、发病机制、临床表现及预防措施；掌握克山病与大骨节病；了解水俣病和宣威肺癌的环境流行病学特征。
}

\section{概述}

\begin{defbox}
\textbf{生物地球化学性疾病（Biogeochemical Disease）}：由于地壳表面化学元素分布的不均匀性，使某些地区的水和/或土壤中某些元素过多或过少，当地居民通过饮水、食物等途径摄入这些元素过多或过少，而引起的特异性疾病。

\textbf{地方病（Endemic Disease）}：由于自然因素和社会因素的影响，在某一地区的人群中发生，不需自外地输入，并呈地方性流行特点的某种疾病。
\end{defbox}

\textbf{流行特征：}①\imp{明显的地区性分布}；②\imp{与环境中元素水平相关}。\textbf{影响因素：}营养条件、生活习惯、多种元素的联合作用。

\textbf{我国主要防治的生物地球化学性疾病（5种）：}碘缺乏病、地方性氟中毒、地方性砷中毒、克山病、大骨节病。

\textbf{综合控制措施：}①限制摄入（改换水源除氟/砷、改良炉灶、改变生活方式）；②适量补充（食盐加碘——最安全有效、口服亚硒酸钠补硒）；③组织措施（建立防治机构、监测病情、动员群众参与）。

\section{碘与健康——碘缺乏病}

\begin{defbox}
\textbf{碘缺乏病（IDD）}：从胚胎发育至成人期由于碘摄入量不足而引起的一系列病症。包括地方性甲状腺肿、地方性克汀病（呆小症）、地方性亚临床克汀病、碘缺乏导致的流产、早产、死产、先天畸形等。

\textbf{地方性克汀病（呆小症）}：由于外环境较严重缺碘引起的以脑发育障碍和体格发育落后为主要特征的地方病。主要表现为较严重的智力障碍、聋哑、神经运动功能障碍、体格发育落后等——\imp{"呆、小、聋、哑、瘫"}。

\textbf{剂量参考：}最低生理需要量75 $\mu$g/d，供给量150 $\mu$g/d，孕妇250 $\mu$g/d；摄入碘低于40 $\mu$g/d、水中含碘量低于10 $\mu$g/L时，碘缺乏病流行。
\end{defbox}

\subsection{流行特征}

\begin{itemize}[leftmargin=*]
    \item 地区分布：\imp{山区 $>$ 丘陵 $>$ 平原 $>$ 沿海}；内陆 $>$ 沿海
    \item 人群分布：任何年龄均可发病，女性早于男性
    \item 时间趋势：补碘干预后可迅速改变流行状况
\end{itemize}

\subsection{发病机制}

碘缺乏 → 甲状腺激素合成不足 → TSH反馈性升高 → 甲状腺代偿性增生 → \imp{地方性甲状腺肿}。重度缺碘影响胎儿脑发育 → \imp{地方性克汀病}（呆小症）。

\subsection{预防措施}

\imp{碘盐}（最经济有效）；碘油（重点人群）；富碘食物（海带、紫菜等）。

\section{氟与健康——地方性氟中毒}

\begin{defbox}
\textbf{地方性氟中毒（Endemic Fluorosis）}：由于一定地区的环境中氟元素过多，导致生活在 该环境中的居民经饮水、食物和空气等途径长期摄入过量氟，引起以\keyword{氟骨症和氟斑牙}为主要特征的慢性全身性疾病。人体摄入总氟量超过4 mg/d即可引起慢性氟中毒。
\end{defbox}

\subsection{发病机制详解}

\begin{enumerate}[leftmargin=*]
    \item \imp{对骨组织的影响}：羟基磷灰石→氟磷灰石→难溶性氟化钙（CaF$_2$）。氟对成骨细胞的作用呈双向性：长期小剂量氟促进骨形成，大剂量诱导细胞凋亡引起骨质疏松。双向作用使氟中毒时出现骨质硬化、骨质疏松或两者同时并存
    \item \imp{对钙磷代谢的影响}：过量氟消耗大量钙→血钙降低→刺激甲状旁腺分泌激素增多→抑制肾小管对磷的重吸收→磷排出增多→磷代谢紊乱。血钙减少和甲状旁腺激素增加又刺激钙从骨组织中不断释放入血，造成骨质脱钙或溶骨，临床上可表现为骨质疏松、骨软化甚至骨骼变形
    \item \imp{抑制酶活性}：氟可与某些酶结构中的金属离子形成复合物，或与带正电的赖氨酸和精氨酸基团、磷蛋白等结合，改变酶结构、抑制酶活性
    \item \imp{对牙齿的影响——氟斑牙}：氟化钙沉积于发育中的牙组织，导致色素沉着、易碎缺损
    \item \imp{非骨相损害}：以中枢神经系统损害最为常见；损伤神经受体；使骨骼肌纤维萎缩；损伤肾小管；诱导肝细胞凋亡；使血管壁钙化硬化；使C细胞功能紊乱；破坏睾丸细胞结构导致生殖功能下降
\end{enumerate}

\subsection{氟的双重健康效应}

\begin{itemize}[leftmargin=*]
    \item \imp{适量}：预防龋齿，促进骨骼钙化
    \item \imp{过量}：氟斑牙、氟骨症
    \item \imp{缺乏}：龋齿发病率增加
\end{itemize}

\subsection{病区类型与预防}

\begin{table}[htbp]
\centering
\caption{地方性氟中毒病区类型及预防措施}
\begin{tabularx}{\textwidth}{lX}
\hline
\textbf{病区类型} & \textbf{预防措施} \\
\hline
\imp{饮水型}（分布最广） & 改换水源（改用低氟水源）；饮水除氟 \\
\imp{燃煤污染型} & 改良炉灶；减少食物氟污染；不用高氟劣质煤 \\
\imp{饮砖茶型} & 研制低氟砖茶；降低砖茶氟含量 \\
\hline
\end{tabularx}
\end{table}

\subsection{人群分布}

年龄：氟斑牙——生长发育中的恒牙（7--8岁前）；氟骨症——成年人。性别：女性以骨软化为主，男性以骨硬化为主。居住时间越长病情越重。

\section{砷与健康——地方性砷中毒}

\begin{defbox}
\textbf{地方性砷中毒（Endemic Arsenism）}：长期从饮用水、室内煤烟、食物等环境介质中摄入过量砷引起以\keyword{皮肤损害}（色素沉着、脱失、角化）为主要特征，伴有多系统、多脏器受损的慢性全身性疾病。砷为IARC确认的\imp{1类致癌物}。

\textbf{砷的代谢：}
\begin{itemize}[leftmargin=*]
    \item 吸收：呼吸道（三价砷，以颗粒物为载体沉积肺组织）；消化道（吸收容易度：五价砷$>$三价砷，无机砷$>$有机砷）；皮肤黏膜（贮存于角蛋白中）
    \item 分布蓄积：三价砷极易与巯基结合→分布于肝、肾、脑等实质性脏器；五价砷取代骨组织中磷灰石的磷酸盐→蓄积于骨组织；三价砷易与角蛋白结合→蓄积于皮肤、指（趾）甲、毛发。毛发砷含量为早期、敏感的内暴露生物标志
    \item 排泄：肾脏是主要排泄器官，尿砷测定可灵敏反映机体砷内暴露水平
\end{itemize}

\textbf{发病机制：}①抑制酶活性——三价砷与酶蛋白双巯基或羧基结合形成稳定络合物，五价砷与酶结合物不稳定可自行水解恢复活性；砷对酪氨酸酶呈双相反应——小剂量产生大量黑色素，大量蓄积致皮肤色素脱失；②导致细胞凋亡；③致癌——直接导致DNA损伤、影响损伤修复、基因表达异常、DNA甲基化反应、产生活性氧引起氧化应激。

\textbf{病区类型：}饮水型（台湾"乌脚病"、新疆、内蒙古）和燃煤污染型（我国特有，分布于贵州、四川、陕西，陕西为"氟砷联合污染型"）。
\end{defbox}

\section{硒与健康}

\subsection{克山病}

\begin{defbox}
\textbf{克山病（Keshan Disease）}：与环境低硒有关的地方性心肌病。1935年首先发现于黑龙江省克山县。
\end{defbox}

\textbf{病因：}环境硒缺乏 + 柯萨奇病毒感染等多因素。\textbf{病理：}心肌变性、坏死、纤维化。\textbf{预防：}口服\imp{亚硒酸钠}补硒，改善膳食营养。

\subsection{大骨节病}

\textbf{大骨节病（Kaschin-Beck Disease）}：与环境低硒有关的慢性骨关节病。以骺板软骨和关节软骨变性坏死为特征，导致关节增粗、短指畸形、身材矮小。

\section{环境污染性疾病}

\begin{defbox}
\textbf{环境污染性致病因素}：能污染环境，使环境质量恶化，而直接或间接使人患病的环境污染因素。

\textbf{环境污染性疾病}：由环境污染因素在暴露人群中引发的疾病。

\textbf{公害病}：严重的地区性环境污染性疾病，由政府认定，须经严格鉴定和国家法律的正式认可，有\imp{医学和法律的双重含义}，病人享受国家补贴。
\end{defbox}

\textbf{环境污染性疾病的特点：}
\begin{enumerate}[leftmargin=*]
    \item 环境污染区域内的人群不分年龄、性别都可能发病
    \item 发病者均出现与暴露污染物相关的相同症状和体征
    \item 除急性危害外，大多具有低浓度长期暴露、陆续发病的特点
    \item 往往缺乏健康危害的早期诊断指标
    \item 预防的关键在于消除环境污染性致病因素、加强对易感人群和亚临床阶段人群的保护
\end{enumerate}

\subsection{慢性甲基汞中毒（水俣病）}

\begin{keybox}
\textbf{水俣病——慢性甲基汞中毒：}
\begin{itemize}[leftmargin=*]
    \item 污染来源：含汞工业废水→无机汞在底泥中经微生物作用甲基化→甲基汞经食物链生物放大（鱼贝类富集百万倍以上）→人食鱼中毒
    \item 发病机制：甲基汞主要经消化道摄入，具有脂溶性、原形蓄积和高神经毒特性。甲基汞与胃酸作用生成氯化甲基汞，吸收入血后与血红蛋白中的巯基结合，随血液分布到脑、肝、肾脏和其他组织。甲基汞能通过血脑屏障进入脑细胞，还能透过胎盘进入胎儿脑中。脑细胞富含类脂质，甲基汞对此有很高亲和力，易蓄积在脑细胞内，造成\imp{进行性不可逆性}的病理损害
    \item 临床表现：感觉障碍（口周和肢端麻木）、共济失调、视野缩小、听力障碍。先天性水俣病：脑性瘫痪——甲基汞通过胎盘导致胎儿中枢神经系统发育障碍
\end{itemize}
\end{keybox}

\subsection{宣威肺癌与军团菌病}

\subsection{军团病}

\begin{keybox}
\textbf{军团病（Legionnaires Disease, LD）}：由\imp{嗜肺军团菌（Legionella Pneumophila, LP）}引起的一种以肺炎为主的全身性疾病，以肺部感染伴全身多系统损伤为主要表现，也可表现为无肺炎的急性自限性流感样疾病。

\textbf{流行病学特征：}
\begin{itemize}[leftmargin=*]
    \item \textbf{传染源}：人工水环境，包括冷热水管道系统、集中空调冷却塔循环水、淋浴设施、游泳池、喷泉、空气加湿器
    \item \textbf{传播途径}：吸入被军团菌污染的气溶胶
    \item \textbf{时间分布}：多发于夏秋季
    \item \textbf{人群分布}：中老年男性和烟酒嗜好者多见
\end{itemize}

\textbf{诊断依据：}
\begin{enumerate}[leftmargin=*]
    \item 流行病学资料：夏秋发病，环境中有建筑施工、空调系统、淋浴喷头设施
    \item 临床资料：临床表现+胸片+非特异性实验室检查（无特异性、多样性）；特异性检查——呼吸道分泌物和肺组织标本培养、血清/呼吸道分泌物抗体荧光抗体法检测、痰/尿液抗原ELISA检测
    \item 环境检测：对室内空气、冷凝塔水等环境样品进行细菌培养
\end{enumerate}
\end{keybox}

\begin{tipbox}
\textbf{宣威女性肺癌高发：}室内燃煤产生PAHs等致癌物→女性暴露更多→肺癌高发。改造炉灶后发病率显著下降。
\end{tipbox}

\section{复习题}

\begin{enumerate}[leftmargin=*, itemsep=8pt]
    \item \textbf{【名词解释】}生物地球化学性疾病；碘缺乏病（IDD）
    \item \textbf{【名词解释】}地方性氟中毒；氟斑牙；氟骨症
    \item \textbf{【名词解释】}克山病；地方性砷中毒
    \item \textbf{【简答题】}简述生物地球化学性疾病的流行特征。
    \item \textbf{【简答题】}简述碘缺乏病的流行特征及预防措施。
    \item \textbf{【简答题】}简述地方性氟中毒的病区类型及预防措施。
    \item \textbf{【简答题】}简述氟的生理作用及双重健康效应。
    \item \textbf{【简答题】}简述慢性甲基汞中毒（水俣病）的病因及发病机制。
    \item \textbf{【简答题】}简述克山病的病因、病理特征及预防措施。
    \item \textbf{【简答题】}简述宣威肺癌高发的环境流行病学研究及其启示。
\end{enumerate}

\subsection*{参考答案要点}

\begin{enumerate}[leftmargin=*, itemsep=5pt]
    \item \textbf{生物地球化学性疾病}：因地球化学元素分布不均，摄入元素过多或过少引起的特异性疾病。\textbf{IDD}：从胚胎至成人期因碘摄入不足引起的一系列病症。

    \item \textbf{地方性氟中毒}：长期摄入过量氟引起的以氟骨症和氟斑牙为特征的慢性全身性疾病。

    \item \textbf{克山病}：与环境低硒有关的地方性心肌病。\textbf{地方性砷中毒}：长期摄入过量砷引起以皮肤损害为特征的慢性中毒。

    \item 两大特征：明显地区性分布；与环境中元素水平相关。影响因素：营养条件、生活习惯、多元素联合作用。

    \item 分布：山区$>$丘陵$>$平原$>$沿海，内陆$>$沿海。预防：碘盐、碘油、富碘食物。

    \item 三类病区：饮水型（改换水源+除氟）、燃煤污染型（改良炉灶）、饮砖茶型（低氟砖茶）。

    \item 适量：预防龋齿、促进骨骼钙化。过量：氟斑牙、氟骨症。缺乏：龋齿增多。呈U型剂量-效应关系。

    \item 来源：含汞废水→甲基化→食物链放大→人食鱼中毒。机制：甲基汞穿过血脑和胎盘屏障，损害中枢神经系统。

    \item 病因：低硒+柯萨奇病毒感染。病理：心肌变性坏死纤维化。预防：口服亚硒酸钠、改善膳食。

    \item 病因：室内燃煤产生PAHs。特点：女性发病率高于男性。启示：流行病学+干预研究整合范式。
\end{enumerate}

% ----- PPT参考题 -----
\subsection*{参考题（PPT课件补充——课程目标自测）}

\begin{enumerate}[leftmargin=*, itemsep=10pt]
    \item \textbf{【简答题】}什么是生物地球化学性疾病？其核心特征是什么？
    \item \textbf{【简答题】}生物地球化学性疾病有哪些主要流行特征？
    \item \textbf{【简答题】}我国主要防治的生物地球化学性疾病有哪些？（至少列出5种）
    \item \textbf{【简答题】}生物地球化学性疾病的流行受哪些影响因素的影响？举例说明。
    \item \textbf{【简答题】}简述生物地球化学性疾病的综合控制措施。
\end{enumerate}

\subsubsection*{参考答案}
\begin{enumerate}[leftmargin=*, itemsep=5pt]
    \item \textbf{定义}：由于地壳表面化学元素分布的不均匀性，使某些地区的水和/或土壤中某些元素过多或过少，当地居民通过饮水、食物等途径摄入这些元素过多或过少，而引起的特异性疾病。\textbf{核心特征}：呈明显地区性分布，与特定地理环境中某元素水平密切相关。
    \item 两大特征：①明显的地区性分布——山区、内陆等特定地理单元高发；②与环境中元素水平密切相关——摄入量过多或过少均可致病。
    \item 碘缺乏病（地方性甲状腺肿、地方性克汀病）、地方性氟中毒（饮水型/燃煤污染型/饮茶型）、地方性砷中毒（饮水型/燃煤污染型）、克山病（低硒相关地方性心肌病）、大骨节病（低硒相关骨关节病）。
    \item 三大影响因素：①营养条件——蛋白质充足可拮抗氟、砷毒性，维生素C促进氟排泄；②生活习惯——如燃煤烘烤食物、饮用砖茶等；③多种元素联合作用——如低碘+低硒协同加重碘缺乏病，高氟+低碘使氟中毒提前出现。
    \item 控制措施：①限制摄入——改换水源除氟/砷、改良炉灶、改变生活方式；②适量补充——食盐加碘（最安全有效）、口服亚硒酸钠补硒；③组织措施——建立防治机构、监测病情、动员群众参与。
\end{enumerate}

\newpage

% =====================================================================
%  CHAPTER 9-10: 公共场所卫生 & 环境质量评价
% =====================================================================
% =====================================================================
%  CHAPTER 9: 家用化学品卫生
% =====================================================================
\chapter{家用化学品卫生}
\setcounter{chapter}{9}
\chapterintro{
了解化妆品的定义、分类及对健康的不良影响；掌握化妆品对皮肤的损害类型；了解化妆品微生物污染的危害。
}

\section{化妆品概述}

\begin{defbox}
\textbf{化妆品（Cosmetic）}：以涂擦、喷洒或其他类似方法，散布于人体任何表面（皮肤、毛发、指甲、口唇等）、牙齿及口腔黏膜，以清洁、保护、美化、修饰及保持其处于良好状态为目的的产品。
\end{defbox}

\subsection{化妆品分类}

\begin{enumerate}[leftmargin=*]
    \item \imp{一般用途化妆品}：护肤类、益发类、美容修饰类、芳香类、口腔卫生用品
    \item \imp{特殊用途化妆品}：用于育发、染发、烫发、脱毛、除臭、祛斑和防晒的化妆品。常需加入某些活性物质以弥补体表局部缺陷，其中有些有一定副作用而在化妆品中被限制使用
\end{enumerate}

\section{化妆品对健康的不良影响}

\subsection{化妆品对皮肤的不良影响}

\begin{keybox}
\textbf{1. 刺激性接触性皮炎（ICD）}：\imp{最常见}的化妆品皮肤损害，因反复接触引起。无免疫记忆，首次接触即可发生。

\textbf{2. 变应性接触性皮炎}：接触化妆品数天后局部出现皮炎症状，属于迟发型变态反应，第二次接触后发生。

\textbf{3. 化妆品光感性皮炎}：化妆品中某些成分在日光（紫外线）作用下引起的光毒性或光变应性反应。包括光变应性接触性皮炎（PCD）和光毒性皮炎。

\textbf{4. 化妆品痤疮}：由矿物油制成的基质引起，油性皮肤者多见。诊断需重点了解化妆品接触史和施用部位。

\textbf{5. 化妆品皮肤色素异常}：化妆品成分导致的色素沉着或色素减退。
\end{keybox}

\subsection{化妆品微生物污染危害}

\begin{enumerate}[leftmargin=*]
    \item \imp{一次污染}：化妆品原料本身有微生物，或在生产过程中受污染
    \item \imp{二次污染}：化妆品启封后，使用或存放过程中受到污染，包括手部接触带入微生物或空气中微生物落入
\end{enumerate}

微生物污染可引起化妆品腐败变质，其在代谢过程中产生的毒素作为变应原或刺激原对施用部位产生致敏或刺激作用。此外，化妆品被致病菌污染可能引起局部甚至全身感染。

\subsection{化妆品所含化学物质的毒性作用}

特殊用途化妆品中常需加入某些活性物质（如染发剂中的对苯二胺、烫发剂中的巯基乙酸等），这些化学物质可能通过皮肤吸收产生全身毒性作用，如致突变、致癌、致畸等。

\section{复习题}

\begin{enumerate}[leftmargin=*, itemsep=8pt]
    \item \textbf{【名词解释】}化妆品；ICD
    \item \textbf{【简答题】}简述化妆品的分类。
    \item \textbf{【简答题】}简述化妆品对皮肤的主要不良影响类型。
    \item \textbf{【简答题】}简述化妆品微生物污染的一次污染与二次污染的区别。
\end{enumerate}

\subsection*{参考答案要点}

\begin{enumerate}[leftmargin=*, itemsep=5pt]
    \item \textbf{化妆品}：以涂擦、喷洒或其他类似方法散布于人体表面、牙齿及口腔黏膜，以清洁、保护、美化、修饰及保持良好状态为目的的产品。\textbf{ICD}：刺激性接触性皮炎，化妆品最常见的皮肤损害。
    \item 一般用途（护肤、益发、美容修饰、芳香、口腔卫生）和特殊用途（育发、染发、烫发、脱毛、除臭、祛斑、防晒）。
    \item 五种类型：刺激性接触性皮炎（ICD，最常见）、变应性接触性皮炎、化妆品光感性皮炎（PCD和光毒性皮炎）、化妆品痤疮、化妆品皮肤色素异常。
    \item 一次污染：原料本身有微生物或生产过程中受污染。二次污染：启封后使用或存放过程中受污染（手部接触或空气中微生物落入）。
\end{enumerate}

\newpage

% =====================================================================
%  CHAPTER 10: 公共场所卫生
% =====================================================================
\chapter{公共场所卫生}
\setcounter{chapter}{10}
\chapterintro{
了解公共场所的概念、分类及卫生学特点；掌握公共场所的基本卫生要求；了解公共场所的卫生管理与监督内容。
}

\section{公共场所概述}

\begin{defbox}
\textbf{公共场所（Public Place）}：根据公众生活活动和社会活动的需要，人工建成的具有多种服务功能的封闭式和开放式的公共建筑设施。
\end{defbox}

\subsection{卫生学特点}

\begin{enumerate}[leftmargin=*]
    \item \imp{人群密集，流动性大}——易造成疾病传播
    \item \imp{设备及物品易被污染}——公共用品反复使用
    \item \imp{涉及面广}——各类人群均可能接触
    \item \imp{从业人员素质参差不齐}——管理水平差异大
\end{enumerate}

\subsection{公共场所的分类}

\begin{multicols}{2}
\begin{enumerate}[leftmargin=*]
    \item 住宿与交际场所（宾馆等）
    \item 洗浴与美容场所
    \item 文化娱乐场所
    \item 体育与游乐场所
    \item 文化交流场所（图书馆等）
    \item 购物场所（商场等）
    \item 就诊与交通场所
\end{enumerate}
\end{multicols}

\section{公共场所的卫生要求与监督}

\textbf{基本卫生要求：}空气质量达标、微小气候适宜、采光照明良好、噪声符合标准、公用物品清洁、卫生设施齐全。

\textbf{经营企业自身卫生管理：}配备卫生管理人员和制度；从业人员卫生培训；健康检查；对顾客卫生宣教。

\textbf{卫生机构管理：}从业人员培训体检；发放卫生许可证；健康教育；经常性卫生监督。

\begin{defbox}
\textbf{经常性卫生监督}：对公共场所日常卫生状况进行的定期巡回监督检查，确保持续符合卫生要求。
\end{defbox}

\section{复习题}

\begin{enumerate}[leftmargin=*, itemsep=8pt]
    \item \textbf{【名词解释】}公共场所；经常性卫生监督
    \item \textbf{【简答题】}简述公共场所的卫生学特点。
    \item \textbf{【简答题】}简述公共场所经营企业自身卫生管理的主要内容。
\end{enumerate}

\subsection*{参考答案要点}

\begin{enumerate}[leftmargin=*, itemsep=5pt]
    \item \textbf{公共场所}：根据公众生活和社会活动需要人工建成的具有多种服务功能的公共建筑设施。\textbf{经常性卫生监督}：对公共场所日常卫生状况的定期巡回监督检查。

    \item 四大特点：人群密集流动性大、设备物品易被污染、涉及面广、从业人员素质参差不齐。

    \item 配备卫生管理人员和制度；组织从业人员学习卫生知识技能；组织健康检查；开展对顾客卫生宣教。
\end{enumerate}

\newpage

% =====================================================================
%  CHAPTER 10: 环境质量评价
% =====================================================================
\chapter{环境质量评价}
\setcounter{chapter}{11}
\chapterintro{
掌握环境质量评价的概念、目的和种类；掌握环境质量指数法的基本原理；了解主要大气质量评价指数（AQI等）和水质评价指数的特点；了解环境影响评价的内容和工作程序。
}

\section{概述}

\begin{defbox}
\textbf{环境质量评价（Environmental Quality Assessment）}：按照一定的评价标准和方法，对一定区域范围内的环境质量进行定性和定量的评定。
\end{defbox}

\textbf{目的：}了解区域环境质量现状和变化趋势；为环境规划管理提供依据；为制定环境标准法规提供参考；评价污染对人群健康的影响。

\textbf{主要内容：}①\imp{污染源评价}——确定主要污染源和污染物；②\imp{环境质量评价}——评价各要素质量；③\imp{环境效应评价}——评价对生态和健康的影响。

\section{环境质量评价方法}

\subsection{环境质量指数法}

\begin{keybox}
\textbf{环境质量指数（EQI）}的计算方法：

\textbf{1. 比值法：}分指数 $I_i = C_i / S_i$（$C_i$为实测浓度，$S_i$为标准值）

\textbf{2. 评分法：}根据污染物浓度分级评分

\textbf{多项分指数合并方法：}简单叠加、算术平均、加权平均、几何平均、Nemerow法（兼顾最大值和平均值）。

\textbf{注意：}需对分指数进行\imp{方向性纠正}：
\begin{itemize}[leftmargin=*]
    \item \textbf{正向指标}（如Cd）：实测值越高，分指数越大，污染越严重
    \item \textbf{逆向指标}（如DO）：实测值越低，分指数越大，污染越严重
    \item \textbf{双向指标}（如pH）：实测值离中心值越远，分指数越大，污染越严重
\end{itemize}
\end{keybox}

\subsection{大气质量评价指数}

\begin{table}[htbp]
\centering
\caption{主要大气质量评价指数及特点}
\begin{tabularx}{\textwidth}{lX}
\hline
\textbf{指数类型} & \textbf{特点} \\
\hline
简单叠加型 & 各分指数相加，简单但信息损失大 \\
比值算术均数型 & 各分指数算术平均 \\
加权算数均数型 & 按重要性加权 \\
I$_1$大气质量指数 & 幂函数模型，考虑S型剂量-反应关系 \\
\imp{AQI（空气质量指数）} & 分段线性函数型，含SO$_2$、NO$_2$、PM$_{10}$、PM$_{2.5}$、CO、O$_3$六项 \\
\hline
\end{tabularx}
\end{table}

\subsection{水质评价指数}

\begin{table}[htbp]
\centering
\caption{主要水质评价指数及特点}
\begin{tabularx}{\textwidth}{lX}
\hline
\textbf{指数} & \textbf{特点} \\
\hline
\imp{Brown水质指数（WQI）} & 评分加权征询法，Delphi法确定权重 \\
\imp{Ross水质指数} & 分级评分叠加型 \\
\imp{综合营养状态指数（TSI）} & 评价富营养化；基准参数：Chl-a、TP、TN、SD、COD$_{Mn}$ \\
\imp{污生指数} & 生物学评价，根据指示生物评价有机污染 \\
\hline
\end{tabularx}
\end{table}

\subsection{土壤环境质量评价}

\begin{defbox}
\textbf{分级污染指数}：按污染物浓度超过背景值的倍数划分污染等级。

\textbf{Nemerow指数}：兼顾平均分指数和最大分指数，$I = \sqrt{\frac{I_{\max}^2 + I_{\text{avg}}^2}{2}}$
\end{defbox}

\section{污染源评价}

\textbf{确定主要污染源和主要污染物：}
\begin{enumerate}[leftmargin=*]
    \item \imp{等标污染负荷法}：$P_i = \frac{Q_i}{C_{0i}} \times 10^{-6}$。将i污染物的排放量稀释到其相应排放标准时所需的介质量，评价各污染源和各污染物的相对危害程度。\textbf{等标污染负荷比}最高的污染物即为最主要的污染物
    \item \imp{排毒系数法}：$F_i = \frac{Q_i}{d_i}$，其中$d_i$为i污染物的评价标准（慢性毒作用阈剂量）。排毒系数表示假设每日排放的i污染物数量长期全部被人们吸入或摄入时，可引起呈现慢性中毒效应的人数
\end{enumerate}

\section{环境影响评价}

\begin{defbox}
\textbf{环境影响评价（EIA）}：对规划和建设项目实施后可能造成的环境影响进行分析、预测和评估，提出预防或减轻不良环境影响的对策和措施。
\end{defbox}

\textbf{我国EIA制度特点：}法定制度；分类管理（报告书/报告表/登记表）；公众参与；资质管理。

\textbf{技术工作程序：}工程分析→环境现状调查评价→环境影响预测评价→环保措施论证→编制报告书。

\section{其他环境要素质量评价}

\subsection{室内环境质量评价}
采用指数综合评价法：
\begin{enumerate}[leftmargin=*]
    \item 确定评价指标，建立分级标准，统一指标方向性
    \item 计算各分指数 $I_i = C_i / S_i$
    \item 计算综合指数
    \item 确定卫生质量等级：一级（很好，$0 \leq Pi < 0.5$）、二级（较好，$0.5 \leq Pi < 1.0$）、三级（较差，$1.0 \leq Pi < 1.5$）、四级（很差，$Pi \geq 1.5$）
\end{enumerate}

\subsection{土壤环境质量评价}
\begin{itemize}[leftmargin=*]
    \item \imp{分级污染指数}：按污染物浓度超过背景值的倍数划分污染等级
    \item \imp{内梅罗（Nemerow）污染指数}：兼顾平均分指数和最大分指数
\end{itemize}

\subsection{生态环境质量评价}
评价生态环境优劣度和动态变化状况。指标包括：生物丰度、生物覆盖指数、水网密度指数、环境质量指数、土地退化指数、生态环境状况指数（EI）。

\section{复习题}

\begin{enumerate}[leftmargin=*, itemsep=8pt]
    \item \textbf{【名词解释】}环境质量评价；环境影响评价（EIA）
    \item \textbf{【名词解释】}环境质量指数（EQI）；等标污染负荷
    \item \textbf{【名词解释】}AQI；综合营养状态指数（TSI）
    \item \textbf{【简答题】}简述环境质量评价的目的和主要内容。
    \item \textbf{【简答题】}简述环境质量指数法的基本原理。
    \item \textbf{【简答题】}AQI日报的主要监测指标有哪些？
    \item \textbf{【简答题】}简述环境影响评价的技术工作程序。
\end{enumerate}

\subsection*{参考答案要点}

\begin{enumerate}[leftmargin=*, itemsep=5pt]
    \item \textbf{环境质量评价}：按一定标准和方法对区域环境质量进行定性和定量评定。\textbf{EIA}：对规划和建设项目可能造成的环境影响进行分析、预测和评估。

    \item \textbf{EQI}：将环境监测数据转化为可比较的无量纲数值的方法。\textbf{等标污染负荷}：污染物排放量与其评价标准值之比。

    \item \textbf{AQI}：空气质量指数（分段线性函数型），包含SO$_2$、NO$_2$、PM$_{10}$、PM$_{2.5}$、CO、O$_3$六项。\textbf{TSI}：以Chl-a、TP、TN、SD、COD$_{Mn}$为基准参数评价富营养化程度。

    \item 目的：了解现状和趋势、为规划管理提供依据、为标准法规提供参考、评价健康影响。内容：污染源评价、环境质量评价、环境效应评价。

    \item 比值法：$I_i=C_i/S_i$。评分法：按浓度分级评分。合并方法：叠加、平均、加权等。需注意方向性统一。

    \item SO$_2$、NO$_2$、PM$_{10}$、PM$_{2.5}$、CO、O$_3$共六项。

    \item 工程分析→环境现状调查评价→环境影响预测评价→环保措施论证→编制报告书。
\end{enumerate}

\newpage

% =====================================================================
%  附录
% =====================================================================
\chapter*{附录一：核心名词解释速记表}
\addcontentsline{toc}{chapter}{附录一：核心名词解释速记表}

\begin{table}[htbp]
\centering
\begin{tabularx}{\textwidth}{lX}
\hline
\textbf{名词} & \textbf{核心释义} \\
\hline
环境卫生学 & 研究自然环境和生活环境与人群健康关系的科学 \\
一次污染物 & 污染源直接排放的未变化污染物 \\
二次污染物 & 一次污染物经物化生作用转化形成的新污染物 \\
生物放大作用 & 重金属/难降解有机物经食物链逐级浓缩放大 \\
人群健康效应谱 & 暴露人群不良反应的五级分布（负荷↑→代偿→亚临床→临床→死亡） \\
健康危险度评价(HRA) & 危害鉴定+暴露评价+剂量反应+风险特征的综合评价 \\
基准 vs 标准 & 科学参考值 vs 法律强制限量值 \\
温室效应 & CO$_2$等温室气体吸收地表热辐射使大气增温 \\
光化学烟雾 & NOx+VOCs+日光紫外线→O$_3$、PAN等强氧化性烟雾 \\
酸雨 & pH$<$5.6的降水，我国以硫酸型为主 \\
逆温 & 大气温度随高度升高而上升，极不利于扩散 \\
PM$_{2.5}$ & 细颗粒物，粒径$\leq$2.5$\mu$m，可入肺泡 \\
水体自净 & 污染物在水体物化生作用下浓度降低的过程 \\
富营养化 & 含P、N污水使藻类大量繁殖（淡水→水华，海湾→赤潮） \\
DO/BOD/COD & 溶解氧/生化需氧量/化学耗氧量 \\
介水传染病 & 通过饮用或接触受病原体污染水引起的传染病 \\
混凝沉淀 & 加入混凝剂使水中悬浮物和胶体聚合沉淀 \\
氯化消毒 & 含氯制剂生成HOCl杀菌 \\
土壤背景值 & 未受污染天然土壤中各元素含量 \\
POPs特性 & 持久性、蓄积性、迁移性、高毒性 \\
生物地球化学性疾病 & 因环境中元素过多或过少引起的特异性疾病 \\
水俣病 & 甲基汞经食物链引起的慢性中毒，损害神经系统 \\
痛痛病 & 镉污染引起的慢性中毒，骨骼病变和肾损害 \\
克山病 & 低硒相关的地方性心肌病 \\
AQI & 空气质量指数，含SO$_2$、NO$_2$、PM$_{10}$、PM$_{2.5}$、CO、O$_3$六项 \\
\hline
\end{tabularx}
\end{table}

\newpage

\vspace*{3cm}
\begin{center}
    {\Large\color{main}\textbf{--- 复习笔记完 ---}}\\[20pt]
    {\normalsize\color{gray!70}
    本笔记依据以下资料整理：\\[5pt]
    《环境卫生学》教学大纲（2026版）\\
    教师授课PPT/PDF课件\\
    往届复习参考资料\\
    教学平台在线题库\\
    \vspace{20pt}
    祝考试顺利！\\[20pt]
    \textit{By greedyCat}
    }
\end{center}

\end{document}
